\documentclass[oneside]{ctexbook}
\usepackage{ctex}
\usepackage[utf8]{inputenc}
\usepackage{amsmath}
%\usepackage{amsthm}
\usepackage{ntheorem}
\usepackage{booktabs}
\usepackage{caption}
\usepackage{listings}
\usepackage[dvipsnames]{xcolor}
\usepackage{xeCJK}
\usepackage{bm}
\usepackage{fancyhdr}
\usepackage{graphicx}
\usepackage{amssymb}
\usepackage{mathrsfs}
\usepackage{titlesec, blindtext, color}
\usepackage{arydshln}
\usepackage{hyperref} 
\usepackage[OT1]{fontenc}
\usepackage{geometry}
\usepackage{comment}
\usepackage{extarrows}
\usepackage{inconsolata}
\usepackage{yhmath}
\usepackage{enumitem}
\usepackage[titletoc]{appendix}
\usepackage{verbatim}

\geometry{a4paper,scale=0.8}

\hypersetup{hidelinks,
	colorlinks=true,
	allcolors=black,
	pdfstartview=Fit,
	breaklinks=true
}

\newcounter{sectionCnt} 


\ctexset{
	 chapter = {
		name = {第,章},
		format = \linespread{1.0}\zihao{1}\bfseries\sffamily\centering,
		nameformat = {\fontsize{30pt}{0pt}\centering\vspace{15pt}},
		titleformat = {\centering},
		number = \chinese{chapter},
		numberformat = \sffamily,
		aftername = ,
		beforeskip = {7pt},
		afterskip = {18pt},
		aftername = {\\},
		pagestyle = plain,
	},
	section={
		%format用于设置章节标题全局格式，作用域为标题和编号
		%字号为小三，字体为黑体，左对齐
		%+号表示在原有格式下附加格式命令
		format+ = \zihao{2} \heiti \raggedright\centering,
		%name用于设置章节编号前后的词语
		%前、后词语用英文状态下,分开
		%如果没有前或后词语可以不填
		name = {,.\,},
		%number用于设置章节编号数字输出格式
		%\chinese输出section编号为中文，\arabic输出为阿拉伯数字
		%number = \Roman{sectionCnt},
		number = \chinese{sectionCnt},
		%beforeskip用于设置章节标题前的垂直间距
		%ex为当前字号下字母x的高度
		%基础高度为1.0ex，可以伸展到1.2ex，也可以收缩到0.8ex
		beforeskip = 1.0ex plus 0.2ex minus .2ex,
		%afterskip用于设置章节标题后的垂直间距
		afterskip = 1.0ex plus 0.2ex minus .2ex,
		%aftername用于控制编号和标题之间的格式
		%\hspace用于增加水平间距
		aftername = \hspace{0pt}
	},
	subsection={
		format+ = \zihao{-2} \kaishu  \raggedright,
		%仅输出subsection编号且为中文
		number = \roman{subsection},
		name = {,.},
		beforeskip = 1.0ex plus 0.2ex minus .2ex,
		afterskip = 1.0ex plus 0.2ex minus .2ex,
		aftername = \hspace{0pt}
	},
	subsubsection={
		%设置对齐方式为居中对齐
		format+ = \zihao{-3} \fangsong,
		%仅输出subsubsection编号，格式为阿拉伯数字，打字机字体
		number = \ttfamily\arabic{subsubsection},
		name = {,.},
		%beforeskip = 1.0ex plus 0.2ex minus .2ex,
		afterskip = 1.0ex plus 0.2ex minus .2ex,
		aftername = \hspace{0pt}
	}
}

\newenvironment{Appendices}[1]
{

	\ctexset{
		chapter = {
			name = {#1},
			format = \linespread{1.0}\zihao{1}\bfseries\sffamily\centering,
			nameformat = {\fontsize{30pt}{0pt}\centering\vspace{15pt}},
			titleformat = {\centering},
			%number = \chinese{chapter},
			number = ,
			numberformat = \sffamily,
			aftername = ,
			beforeskip = {7pt},
			afterskip = {18pt},
			aftername = {\\},
			pagestyle = plain,
		},
		section={
			%format用于设置章节标题全局格式，作用域为标题和编号
			%字号为小三，字体为黑体，左对齐
			%+号表示在原有格式下附加格式命令
			format+ = \zihao{-1} \heiti \raggedright\centering,
			%name用于设置章节编号前后的词语
			%前、后词语用英文状态下,分开
			%如果没有前或后词语可以不填
			name = {,.\,},
			%number用于设置章节编号数字输出格式
			%\chinese输出section编号为中文，\arabic输出为阿拉伯数字
			number = \Roman{section},
			%beforeskip用于设置章节标题前的垂直间距
			%ex为当前字号下字母x的高度
			%基础高度为1.0ex，可以伸展到1.2ex，也可以收缩到0.8ex
			beforeskip = 1.0ex plus 0.2ex minus .2ex,
			%afterskip用于设置章节标题后的垂直间距
			afterskip = 1.0ex plus 0.2ex minus .2ex,
			%aftername用于控制编号和标题之间的格式
			%\hspace用于增加水平间距
			aftername = \hspace{0pt}
		},
		subsection={
			format+ = \zihao{-2} \kaishu  \raggedright,
			%仅输出subsection编号且为中文
			number = \roman{subsection},
			name = {,.},
			beforeskip = 1.0ex plus 0.2ex minus .2ex,
			afterskip = 1.0ex plus 0.2ex minus .2ex,
			aftername = \hspace{0pt}
		},
		subsubsection={
			%设置对齐方式为居中对齐
			format+ = \zihao{-3} \fangsong,
			%仅输出subsubsection编号，格式为阿拉伯数字，打字机字体
			number = \ttfamily\arabic{subsubsection},
			name = {,.},
			%beforeskip = 1.0ex plus 0.2ex minus .2ex,
			afterskip = 1.0ex plus 0.2ex minus .2ex,
			aftername = \hspace{0pt}
		}
	}
}

\fancyhf{}
\fancyhead[L]{
	\begin{minipage}[c]{0.06\textwidth}
		\hyperlink{Index}{\includegraphics[height=7.5mm]{1.jpg}}
	\end{minipage}
	\begin{minipage}[c]{0.4\textwidth}
	线性代数提高笔记
	\end{minipage}}
\fancyhead[R]{
    \begin{minipage}[r]{0.1\textwidth}
	\href{https://qifengggg.github.io/}{奇峰}
    \end{minipage}
}
\cfoot{\hyperlink{Index}{\thepage}}

\pagestyle{fancy}

\fancypagestyle{plain}{
  \fancyhead{}
  \renewcommand{\headrulewidth}{0pt}
}
%\renewcommand{\qedsymbol}{$\blacksquare$}

\renewcommand{\labelenumi}{\roman{enumi}.}
\renewcommand{\labelitemii}{$ \circ $ }
\newcommand{\dis}{\displaystyle}
\newcommand{\Attention}[1]{\textcolor{red}{\bf{#1}}}
\newcommand{\red}{\textcolor{red}}
\newcommand{\sssubsection}[1]{\noindent\textbf{#1}}
\newcommand{\nextline}{\vspace{5pt}\newline}
\newcommand{\dsqrt}{\dis\sqrt}
\newcommand{\pprime}{{\prime\prime}}
\newcommand{\whichis}{\mathop=\limits^\triangle}
\newcommand{\R}{\mathbb R}
\newcommand{\qline}{\underline{~~~~~~~~~~}}

\newcommand{\Section}{\stepcounter{sectionCnt}\section}

\title{\includegraphics[scale=0.6]{1.jpg}\\ \textsf{线性代数提高笔记}}
\author{奇峰}
\date{之前}

\begin{document}

\theoremseparator{}
\newtheorem{def1}{定义}[section]
\newtheorem{theo1}{定理}[section]
\newtheorem{func1}{方法}[section]
\newtheorem{infer1}{推论}[section]
\newenvironment{proof}{\begin{itemize}\item \textbf{证明}
	
	}{\end{itemize}}
%\newenvironment{<环境名称>}[<参数个数>][<首参数默认值>]{<环境前定义>}
%							{<环境后定义>}


\newenvironment{Def}[1][\quad{}]{\begin{def1}\textbf{#1}}{\end{def1}}
\newenvironment{Theo}[1][\quad{}]{\begin{theo1}\textbf{#1}}{\end{theo1}}
\newenvironment{Func}[1][\quad{}]{\begin{func1}\textbf{#1}}{\end{func1}}
\newenvironment{Infer}[1][\quad{}]{\begin{infer1}\textbf{#1}}{\end{infer1}}
\newenvironment{Field}[1][\quad{}]{\noindent\newline\textbf{#1}}{}

\setlength{\parskip}{3pt}
\setenumerate[1]{topsep=2pt}
\setitemize[1]{topsep=2pt}

\frontmatter
\maketitle
\hypertarget{Index}{}
\tableofcontents

\mainmatter

\chapter*{大作文的格式}

\section{行文逻辑}

\begin{itemize}
    \item \textbf{第一段}
    
    描述漫画内容。

    \begin{enumerate}[label = \Alph*.]
        \item 极简主义 - 抓住与漫画主题相关的细节；
        \item 时态 - 现在进行时、一般现在时；
        \item 汉语提示 - 不要直译
    \end{enumerate}
    \item \textbf{第二段}
    
    \textbf{论证某种品质的重要性;}

    阐述现象重要性;
    \item \textbf{第三段}
    
    提出建议，说明怎么做。
\end{itemize}

\newpage

\section{模板}

\subsection{单幅漫画}

A/An encouraging and thought-provoking story 
emerges from the image presented above: \textrm{(I)}
Apparently, the picture does mirror a common phenomenon in our contemporary society.

As a matter of fact, the picture
attempts to convey a crucial message:
\textrm{(I)} does count in \textrm{(I)} .
To put it more explicitly, it is undeniable that 
\textrm{(II)} (为什么需要XXX/XXX的重要性).
In such cases, \textrm{(III)} ought to \textrm{(III)}
(in order to \dots).
By contrast, should \textrm{(IV)} lack such awareness, \textrm{(IV)}
would be more likely to \textrm{(IV)} in the end.
A host of convincing examples can be found to illustrate my viewpoint.

On the whole, according to the analysis mentioned above, we could draw a reasonable 
conclusion that \textrm{(I)} are advised to cultivate/foster a sense of \textrm{(I)}
and put it into practice. Only in this manner can \textrm{(II)}.

%\newpage

\subsection{对比漫画}

The images pose a stark contrast:
in the first picture, \textrm{(I)}.
Conversely, in the second one, \textrm{(II)}.
Apparently, the drawing do reflect distinct attitudes and approaches 
adopted by different people in our contemporary society.

As a matter of fact, the pictures
attempt to convey a crucial message:
\textrm{(I)} does count in \textrm{(I)} .
To put it more explicitly, it is undeniable that 
\textrm{(II)} (为什么需要XXX/XXX的重要性).
In such cases, \textrm{(III)} ought to \textrm{(III)}
(in order to \dots).
By contrast, should \textrm{(IV)} lack such awareness, \textrm{(IV)}
would be more likely to \textrm{(IV)} in the end.
A host of convincing examples can be found to illustrate my viewpoint.

On the whole, according to the analysis mentioned above, we could draw a reasonable 
conclusion that \textrm{(I)} are advised to cultivate/foster a sense of \textrm{(I)}
and put it into practice. Only in this manner can \textrm{(II)}.
\chapter{行列式}

\sssubsection{定义与性质}

$$
    \begin{cases}
        \textrm{三大定义} 
        \begin{cases}
            &\textrm{几何定义};\\
            &\textrm{逆序定义};\\
            &\textrm{展开定义};\\
        \end{cases}
        \\
        \textrm{计算} 
        \begin{cases}
            &\textrm{数值形};\\
            &\textrm{含参形};\\
            &\Attention{\textbf{抽象形}};\\
        \end{cases}
        \\
        \textrm{证明} |A| = 0 
        \begin{cases}
            & |A| = -|A|;\\
            & r(A) < n;\\
            & \textrm{不可逆（常用于反证）};\\
            & AX = 0 \textrm{有非零解};\\
            & \exists \lambda_0 = 0;\\
        \end{cases}
        \\
    \end{cases}
$$ 

\section{求行列式的值}

\begin{itemize}
    \item 用性质消零（展开定义）
    \item 特殊行列式\begin{itemize}
        \item 三角形
        \item 范德蒙德行列式
        \item 分块
    \end{itemize}
    \item 特殊形状的行列式
\end{itemize}

\section{代数余子式}

代数余子式 $ A_{ij} = (-1)^{i+j}M_{ij}. $ 

注意，
\begin{itemize}
    \item 其为 $ (n-1) $ 阶子式;
    \item 其线性组合 $ \sum_i a_iA_ki $ 相当于将行列式第 $ k $ 行元素替换为 $ (a_i); $
    \item $ a_{i1}A_{j1} + a_{i2}A_{j2} + \cdots + a_{in}A_{jn} =
    \begin{cases}D,&i=j\\ 0,i\neq j\end{cases} $ 
    \item 伴随矩阵 $ A^* = \left(A_{ij}\right)^\top\whichis \Attention{|A|\cdot A^{-1}}$ 
\end{itemize}

对数值型矩阵 $ A, $ 求 $ A^* $ 时，可先求 $ |A| $ 并利用 $ (A\, \vdots\, E)\xrightarrow{\textrm{行变换}}
(E\, \vdots\, A^{-1}) $ 
求 $ A^{-1}, $ 再利用公式 $ A^* = |A|A^{-1} $ 求伴随矩阵。

\section{抽象矩阵行列式}

\sssubsection{性质}

\begin{enumerate}
    \item 对已知 $ A = (\alpha_i) $ 的矩阵，可以
    $ \begin{cases}
        \textrm{行列式列消}\\ \textrm{可逆时，}\begin{cases}
            \textrm{矩阵乘法}\\\textrm{相似}
        \end{cases}
    \end{cases} $ 
    \item $ |kA| = k^n|A|, $ 其中 $ k $ 可以是行列式；
    
    $ |A^\top|=A,|AB|=|A||B|=|B||A|=|BA|; $ 
    \item 设 $ A $ 可逆，则 $ A^* $ 可逆，且有

    $ |A^*| = |A|^{n-1};(A^*)^* = |A|^{n-2}A; $ 

    $ |(A^*)^*| = A^{(n-1)^2}; $ 
    \item $ |A| = \prod \lambda_i; tr(A) = \sum a_{ii} = \sum \lambda_i; $ 
    \item $ P^{-1}AP = B \Rightarrow |A| = |B|; $ 
    \item $ aA + bE $ 不可逆 $ \Leftrightarrow |aA + bE| = 0 \Rightarrow \exists \lambda_0 = \dfrac{-b}{a}; $ 
    \item $ AA^\top = A^\top A = E \Rightarrow |A| = \pm 1; $ 
    \item 对可逆矩阵 $ A, |A + \alpha\beta^T| = |A|(1+\beta^TA^{-1}\alpha); $ 
    
    对 $ B = \alpha\beta^T, |\lambda E + \alpha\beta^T| = \lambda^{n-1}(\lambda + \beta^\top\alpha); $ 
    \item 对 $ A_n, $ 若 $ A^2 = A,A\neq E, $ 则有 $ |A| = 0. $ \begin{itemize}
        \item 反证
        
        假设 $ A $ 可逆，则 $ A^2A^{-1} = AA^{-1} = E \Rightarrow A = E, $ 与题设矛盾，故
        $ A $ 不可逆，因此 $ |A| = 0. $ 
        \item 秩
        
        $ A^2 = A \Rightarrow A(A-E) = O \Rightarrow r(A) + r(A-E) \leq n;$
        
        $ A\neq E \Rightarrow r(A-E)>0 \Rightarrow r(A) < n \Rightarrow |A| = 0. $ 
        \item 方程组
        
        $ A^2 = A \Rightarrow A(A-E) = O; A \neq E \Rightarrow AX = O $ 存在非零解，
        因此 $ r(A)<n, $ 即 $ |A|=0. $ 
        \item 特征值
        
        $ A(A-E) = O = O(A-E),A\neq E, $ 因此 $ A $ 存在一特征值 $ \lambda_0 = 0, $ 
        此时 $ |A| = \prod \lambda = 0. $ 
    \end{itemize}
\end{enumerate}

其中，对于第八条，
\begin{enumerate}
    \item 有分块乘法
    \begin{equation*}
        \begin{aligned}
            &\begin{pmatrix}
                E&-\alpha\\0&1
            \end{pmatrix}
            \begin{pmatrix}
                A&\alpha\\-\beta^\top&1
            \end{pmatrix} = 
            \begin{pmatrix}
                A+\alpha\beta^\top & 0 \\ -\beta^T & 1
            \end{pmatrix}&(i)\\
            &\begin{pmatrix}
                E&0\\\beta^\top A^{-1} & 1
            \end{pmatrix}
            \begin{pmatrix}
                A&\alpha\\-\beta^\top&1
            \end{pmatrix} = 
            \begin{pmatrix}
                A & \alpha \\ 0 & 1+\beta^TA^{-1}\alpha
            \end{pmatrix}&(ii)\\
        \end{aligned}
    \end{equation*}
    显然 $ (i),(ii) $ 的两边也相等。此时对二者右边取得行列式，得到$$
    |A + \alpha\beta^T| = |A|(1+\beta^TA^{-1}\alpha)
    $$ 

    考虑一个例子 $ D = \begin{pmatrix}
        x+a_1 & a_2 & \cdots& a_n\\ 
        a_1 & x+a_2 & \cdots& a_n\\ 
        \vdots & \vdots & \ddots & \vdots \\ 
        a_1 & a_2 & \cdots & x+a_n
    \end{pmatrix} $ 
    
    显然其可以被拆分为 $ \begin{pmatrix}
        x&&&\\ &x&& \\ &&\ddots& \\ &&&x
    \end{pmatrix} $ 和 $ \begin{pmatrix}
        a_1 & a_2 & \cdots& a_n\\ 
        a_1 & a_2 & \cdots& a_n\\ 
        \vdots & \vdots & \ddots & \vdots \\ 
        a_1 & a_2 & \cdots & a_n
    \end{pmatrix} $ 
    
    此时可令前者为 $ A, $ 后者为 $ \alpha\beta^\top. $ 
    \item $ |\lambda E + \alpha\beta^\top| = |\lambda E - (-\alpha\beta^\top)|, $ 
    此时视 $ A = \lambda E, $ 
    
    套用前述方法有 $ |\lambda E + \alpha\beta^\top| = 
    |\lambda E|(1+\beta^T(\lambda E)^{-1}\alpha) = \lambda^{n}(1+\frac{\beta^\top\alpha}{\lambda})
    = \lambda^{n-1}(\lambda + \beta^\top\alpha). $ 
\end{enumerate}

注意，若有 $ AB=O,B = (\beta_i), $ 则考虑
\begin{itemize}
    \item $ r(A)+r(B)\leq n; $ 
    \item $ B $ 的列向量为 $ AX = O $ 的一组解；
    \item $ A\beta_i = O = O \beta_i, $ 此时 $ B $ 的列向量为 $ \lambda_A = 0 $ 的特征向量。
\end{itemize}

\sssubsection{计算抽象行列式}

计算抽象的行列式时，若其有法则，如 $ A^{-1},A^\top,A^*, $ 则利用其法则，
若其无法则，则利用 $ E = AA^{-1} $ 或 $ E = AA^\top. $

对求 $ |A+B| $ 的情况，由于无法则，考虑添加 $ E, $ 此时“一前一后，前者前，后者后。”对于

$$
    |A + B| = |E_1A + BE_2|,\ \ E_1,E_2 \textrm{都是单位矩阵}
$$ 

有\begin{itemize}
    \item $ E_1 $ 在 $ A $ 前，$ E_2 $ 在 $ B $ 后，此谓一前一后；
    \item $ E_1 $ 拆时，考虑同在前面的 $ B, $ 即 $ B^{-1}B $ 或 $ BB^{-1}, $ 此谓前者前；
    \item $ E_2 $ 拆时，考虑同在后面的 $ A, $ 即 $ A^{-1}A $ 或 $ AA^{-1}, $ 此谓后者后；
\end{itemize}

\sssubsection{由行列式值求参数}

需要加减·消元·出公因式。如
$$
    \begin{pmatrix}
        \lambda - 3 & 1 & -1 \\ 
        1 &\lambda -5 & 1 \\ 
        -1 & 1 & \lambda - 3
    \end{pmatrix}
    \xlongrightarrow{\textrm{第三列}-1\textrm{倍加到第一列}}
    \begin{pmatrix}
        \lambda - 2 & 0 & 2 - \lambda \\ 
        1 &\lambda -5 & 1 \\ 
        -1 & 1 & \lambda - 3
    \end{pmatrix}
$$ 

此时可以从第一列提出系数 $ \lambda - 2. $ 

\section{求行列式方程的根}

考察的方式可能为
\begin{itemize}
    \item 讨论根的个数 （行列式和最高次方）
    \item 根与系数的关系
    
    对 $ f(x) = D = \sum_{i = 0}^n a_i x^{n-i} $ 及其根 $ x_i $ 
    有 $ \begin{cases}
        \sum x_i = -\dfrac{a_1}{a_0} \\ 
        \prod (-x_i) = \dfrac{a_n}{a_0}
    \end{cases} $ 
\end{itemize}

\sssubsection{矩阵加边}

对行列式 $ D_n, $ 显然$$
    |D| = \left|
    \begin{array}{c:ccc}
        1&*&\cdots&*\\ \hdashline
        0& & & \\
        \vdots& &D& \\
        0& & & \\
    \end{array}
    \right|
$$ 

其中，为 $ * $ 的部分可以为任意值，因此可以按题面设置易于计算的值。另外，加入的一行
具体在哪一行都可以。

也可以通过加边构造行列式，使得其虽然不等于原式，但能通过性质辅助计算。如证明

$$
    D = \begin{vmatrix}
        1&1&1 &1 \\
        a_1^1&a_2^1 &a_3^1 &a_4^1 \\
        a_1^2&a_2^2 &a_3^2 &a_4^2 \\
        a_1^4&a_2^4 &a_3^4 &a_4^4 \\
    \end{vmatrix} = \sum a_i \prod (a_i - a_j)
$$ 

可以将其变为$$
    f(x) = \begin{vmatrix}
        x^0&1&1&1 &1 \\
        x^1&a_1^1&a_2^1 &a_3^1 &a_4^1 \\
        x^2&a_1^2&a_2^2 &a_3^2 &a_4^2 \\
        x^3&a_1^3&a_2^3 &a_3^3 &a_4^3 \\
        x^4&a_1^4&a_2^4 &a_3^4 &a_4^4 \\
    \end{vmatrix} \begin{matrix}
        \\ \\ \\ \leftarrow \textrm{增补行} \\ \\
    \end{matrix}
$$ 

并按照第一列展开，此时 $ x^3 $ 的系数$ A_3 $ 正好为 $ -|D|, $ 而
$ x^4 $ 的系数$ A_4 $ 为范德蒙德行列式，值为 $ \prod (x_i-x_j), $ 
由前述结论，$ \sum a_i = -\dfrac{A_3}{A_4}, $ 整理即得到结论。

\chapter{一元函数的微分学}

\section{导数}

\subsection{导数的概念}

\begin{Def}[导数]

    设函数 $ y=f(x) $ 在点 $ x_{0} $ 的一邻域内有定义，自变量 $ x $ 在$ x_0 $ 处有增量 $ \Delta x $，因变量 $ y $ 
    相应地有增量 $ \Delta y=f(x_0+\Delta x)-f(x_0) $ ，若 
    $ {\displaystyle\lim_{\Delta x\rightarrow 0}}\frac{\Delta y}{\Delta x}=
    {\displaystyle\lim_{\Delta x\rightarrow 0}}\frac{f(x+\Delta x)-f(x)}{\Delta x} $ 存在，则称此
    极限值为函数$ y=f(x) $ 在 $ x=x_0 $ 点处的导数，也称微商，记作$$
        f'(x)\textrm{或}\left.y'\right|_{x=x_0}\textrm{或}\left.\frac{\mathrm{d}f(x)}{\mathrm{d}x}\right|_{x=x_0}
    $$ 
    称函数 $ y=f(x) $ 在 $ x=x_0 $ 处可导。若上述极限不存在，则称 $ f(x) $ 在 $ x=x_0 $ 处不可导。 
\end{Def}

\begin{Field}[导数定义的等价变形]

    导数定义可以等价变形为
    $$
        f'(x)={\displaystyle\lim_{\Delta x\rightarrow 0}}\frac{f(x+\Delta x)-f(x)}{\Delta x}=
        {\displaystyle\lim_{x\rightarrow x_0}} \frac{f(x+x_0)-f(x)}{x-x_0}
    $$ 
\end{Field}
前者适合用于求导函数。

注意：
\begin{itemize}
    \item 定义式中分子和分母的增量必须相对应，如$$
        {\displaystyle\lim_{\Delta x\rightarrow 0}\frac{f(x+2\Delta x)-f(x)}{\Delta x}}\neq f'(x)
    $$ 
    \item 在定义式中，分子中两项应当“一动一静”，否则无法推出可导性。注意，$$
        f(x)\textrm{在}x_0\textrm{处可导}\Rightarrow 
        {\displaystyle\lim_{\Delta x\rightarrow 0}}\frac{f(x+\Delta x)-f(x-\Delta x)}{\Delta x}
    $$ 
    $ \Leftarrow $ 反例 - $ f(x)=|x|, x_0 = 0; $ 
    \item 若已知 $ f(x) $ 在 一点 $ x_0 $ 可导，则可能可以利用导数定义。
\end{itemize}

\begin{Def}[左、右导数]

    若$$
        {\displaystyle\lim_{x\rightarrow x_0^-}}\frac{f(x)-f(x_0)}{x-x_0}\mathop{=}\limits^\exists f_-'(x_0)
    $$ 
    则称 $ f(x) $ 在  $ x_0 $ 处左可导， $ f_-'(x_0) $ 称为其左导数；
    若$$
    {\displaystyle\lim_{x\rightarrow x_0^+}}\frac{f(x)-f(x_0)}{x-x_0}\mathop{=}\limits^\exists f_+'(x_0)
$$ 
则称 $ f(x) $ 在  $ x_0 $ 处右可导， $ f_+'(x_0) $ 称为其右导数；
\end{Def}

\begin{Field}[函数在区间上的可导性]

    \begin{itemize}
        \item 若对 $ \forall x\in (a,b) $ 有 $ f(x) $ 在 $ x_0 $ 可导，则称 $ f(x) $ 在 $ (a,b) $ 上可导；
        \item 若对 $ \forall x\in (a,b) $ 有 $ f(x) $ 在 $ x_0 $ 可导，
        且其在 $ a $ 点右可导，在 $ b $ 点左可导，则称 $ f(x) $ 在 $ [a,b] $ 上可导。
    \end{itemize}
\end{Field}

\subsection{函数可导的条件}

\begin{Theo}[函数可导的必要条件]

    可导必连续。
\end{Theo}
$ \Leftarrow $ 反例如下。
\begin{itemize}
    \item $ f(x)=|x|,x_0=0 $ ;
    \item $ f(x)=\sqrt[3]x,x_0=0 $ 此时其左右导数均无穷大。
\end{itemize}

\begin{Theo}[函数可导的充要条件]

    $ f(x) $ 在 $ x=x_0 $ 处可导$ \Leftrightarrow f_-'(x) $ 与 $ f_+'(x) $  均存在且相等。
\end{Theo}

其常用于：\begin{enumerate}
    \item 已知可导性求参数；
    \item （后者）分段函数分段点的可导性。
\end{enumerate}

\subsection{导数的几何意义}

若函数 $ y=f(x) $ 在 $ x_0 $ 处导数 $ f'(x_0) $ \textbf{存在}，则在几何上， $ f'(x_0) $ 表示曲线
$ y=f(x) $ 在点 $ (x_0,f(x_0)) $ 处的切线斜率，即\begin{itemize}
    \item 切线方程$ y-f(x_0)=f'(x_0)(x-x_0) $ 
    \item 法线方程$ y-f(x_0)=-\frac1{f'(x_0)}(x-x_0) $ 
\end{itemize}

注意：
\begin{itemize}
    \item 求切法线的步骤：\begin{enumerate}
        \item 求切点坐标；
        \item 求斜率；
        \item 写方程。
    \end{enumerate}
    \item 可用于处理\begin{itemize}
        \item 显函数 $ y=f(x) $ ;
        \item 隐函数 $ F(x,y)=0 $ ;
        \item 参数方程、极坐标方程，极坐标时可以代$ r $ 入$ \theta. $
    \end{itemize}
\end{itemize}

\subsection{导数的运算}

以下是一些基本初等函数的导数公式：

\begin{equation*}
    \begin{array}{ll}
        (c)'=0&(x^a)'=ax^{a-1}\\
        (\sin x)'=\cos x&(\cos x)'=-\sin x\\ 
        (\tan x)'=\sec^2 x&(\cot x)'=-\csc^2 x\\ 
        (\sec x)'=\tan x\sec x&(\csc x)'=-\cot x\csc x\\ 
        (\arcsin x)'=\frac{1}{\sqrt{(1-x^2)}}&(\arccos x)'=-\frac{1}{\sqrt{(1-x^2)}}\\
        (\arctan x)'=\frac1{1+x^2} & (\textrm{arccot} x)'=-\frac1{1+x^2}\\
        (a^x)'=a^x\ln a&(\log_a x)'=\frac1{x\ln a}\\
    \end{array}
\end{equation*}

注意到对于三角函数，带有“正”字的正弦、正切、正割、反正弦、反正切的导数都不包含负号，
而带有“余”字的余弦、余切、余割、反余弦、反余切的导数都包含负号。

\begin{Field}[导数的四则运算法则]

    设 $ f(x) $ 与 $ g(x) $ 均在 $ x $ 点可导，则
    \begin{enumerate}
        \item $ \left[f(x)\pm g(x)\right]'=f'(x)\pm g'(x) $ ;
        \item $ \left[f(x)g(x)\right]'=f'(x)g(x)+g'(x)f(x) $ ;
        \item $\dis \left[\frac{f(x)}{g(x)}\right]'=\frac{f'(x)g(x)-g'(x)f(x)}{g^2(x)}$，其中要求 $ g(x)\neq 0 $  ;
    \end{enumerate}
\end{Field}

注意，若 $ f(x) $ 可导但是 $ g(x) $ 不可导于 $ x=x_0, $ 则$ f(x)\pm g(x) $ 也不可导于 $ x=x_0. $ 

可以将和、差、积的公式推广至有限多个。对乘法，有
$$
    \left(\prod_{i=1}^nf_i(x)\right)'=\sum_{i=1}^n\left(f'_i(x)\prod_{\begin{subarray}{align}j=1\\j\neq i\end{subarray}}^n f_j(x)\right)
$$ 

\begin{Field}[行列式函数求导]

    行列式函数求导可以按行或按列求导。
    
    如对二阶行列式函数，有
    \begin{equation*}
        \begin{aligned}
            \left[\begin{matrix}
                a_{11}&a_{12}\\a_{21}&a_{22}
            \end{matrix}\right]'&=(a_{11}a_{22}-a_{12}a_{21})'\\
            &=a_{11}a_{12}'+a_{11}'a_{22}-a_{12}a_{21}'-a_{12}'a_{21}\\ 
            &=\begin{bmatrix}
                a_{11}'&a_{12}'\\a_{21}&a_{22}
            \end{bmatrix}+
            \begin{bmatrix}
                a_{11}&a_{12}\\a_{21}'&a_{22}'
            \end{bmatrix}\\
            &=\begin{bmatrix}
                a_{11}'&a_{12}\\a_{21}'&a_{22}
            \end{bmatrix}+
            \begin{bmatrix}
                a_{11}&a_{12}'\\a_{21}&a_{22}'
            \end{bmatrix}\\
        \end{aligned}
    \end{equation*}
\end{Field}

\begin{Field}[复合函数求导]

    设 $ y=f(u) $ ，$ u=\varphi(x) $ ，若 $ \varphi(x) $ 在 $ x $ 处可导，且 $ f(u) $ 在对应的 $ u $ 处可导，
    则复合函数 $ y=f(\varphi(x)) $ 在 $ x $ 处可导，且有$$
        \frac{\mathrm{d}y}{\mathrm{d}x}=\frac{\mathrm{d}y}{\mathrm{d}u}\cdot
        \frac{\mathrm{d}u}{\mathrm{d}x}=f'(\varphi(x))\varphi'(x)
    $$  
    对应地，有$$
        \mathrm{d}y=f'(u)\mathrm{d}u=f'(\varphi(x))\varphi'(x)\mathrm{d}x
    $$ 
\end{Field}

\begin{Field}[隐函数求导数]

    设 $ y=y(x) $ 由方程 $ F(x,y)=0 $ 确定，则$$
        \frac{\mathrm{d}y}{\mathrm{d}x}={\dfrac{\mathrm{d}F}{\mathrm{d}x}}\Big\slash{\dfrac{\mathrm{d}F}{\mathrm{d}y}}
        ,\quad{}\frac{\mathrm{d}^2F}{\mathrm{d}x^2}=\frac{\mathrm{d}}{\mathrm{d}x}(\frac{\mathrm{d}F}{\mathrm{d}x})
    $$ 
\end{Field}

\sssubsection{反函数求导数}

设 $ x = f^{-1}(y) $ 由 $ y = f(x) $ 确定，则
\begin{itemize}
    \item 若 $ f(x) $ 可导且 $ f'(x)\neq 0, $ 则$ \dis \dfrac{\mathrm{d}x}{\mathrm{d}y} = 
    \dfrac{1}{\mathrm{d}y/ \mathrm{d}x} = \dfrac{1}{f'(x)}. $ 
    \item 若 $ f(x) $ 二阶可导且 $ f'(x)\neq 0 $ 则 $ \dfrac{\mathrm{d}^2y}{\mathrm{d}x^2} = 
    -\dfrac{f^\pprime(x)}{[f'(x)]^3}. $ 
\end{itemize}

$ f(x) $ 可导且 $ f'(x)\neq 0 $ 说明 $ f(x) $ 单调且存在反函数。

\sssubsection{参数方程求导数}

设 $ y = f(x) $ 由参数方程 $ \begin{cases}
    x = x(t) \\ y = y(t)
\end{cases} $ 确定。此时，$ t = t(x),y = y(t(x)). $ 

对于其导数，有
\begin{itemize}
    \item 若 $ x(t),y(t) $ 均可导，且$ x'(t)\neq 0, $ 则 $ \dfrac{\mathrm{d}y}{\mathrm{d}x} = 
    \dfrac{\mathrm{d}y/\mathrm{d}t}{\mathrm{d}x/\mathrm{d}t} = \dfrac{y'(t)}{x'(t)}; $ 
    \item 若 $ x(t),y(t) $ 均二阶可导，且$ x'(t)\neq 0, $ 则 $ \dfrac{\mathrm{d}^2y}{\mathrm{d}x^2} = 
    \dfrac{y^\pprime(t)x'(t) - y'(t)x^\pprime(t)}{[x'(t)]^3}. $ 
\end{itemize}


\begin{Field}[常见初等函数的$ n $ 阶导数公式]

    \begin{enumerate}
        \item $ y=e^x \Rightarrow y^{(n)}=e^x; $ 
        \item $ y=a^x(a>0,a\neq1)\Rightarrow y^{(n)}=a^x(\ln a)^n; $ 
        \item $ y=\sin x \Rightarrow y^{(n)}=\sin(x+\frac{n\pi}{2});$
        \item $ y=\cos x \Rightarrow y^{(n)}=\cos(x+\frac{n\pi}{2});$  
        \item $ (\frac{1}{x+a})^{(n)} = \frac{(-1)^{n}n!}{(x+a)^{n+1}}$ 
        \item $ y=\ln(x+a)\Rightarrow y^{(n)}=(\frac{1}{x+a})^{(n-1)} = \frac{(-1)^{n-1}(n-1)!}{(x+a)^{n}} $ 
    \end{enumerate}
\end{Field}


\section{微分}

\subsection{微分的概念}

\begin{Def}[微分]

    设函数 $ y=f(x) $ 在 $ x=x_0 $ 处有增量 $ \Delta x $ 时，若函数增量 $ \Delta y=f(x_0+\Delta x)=f(x_0) $ 
    可以表为 $ \Delta y=A(x_0)\Delta x+o(\Delta x)(\Delta x\rightarrow 0) $ ，其中 $ A(x_0) $ 与 $ \Delta x $ 
    无关，则称 $ f(x) $ 在 $ x=x_0 $ 处可微，$ \Delta y $ 中的主要线性部分（也称\textbf{线性主部}）
    $ A(x_0)\Delta x $ 称为 $ f(x) $ 在 $ x=x_0 $ 处的微分，记为
    $$
        \mathrm{d}y\arrowvert_{x=x_0}\ \textrm{或}\ \mathrm{d}f(x)\arrowvert_{x+0}
    $$ 
\end{Def}

\subsection{可微与导数的的关系}

\begin{Theo}[可微与导数的关系]

    $ f(x) $ 在 $ x=x_0 $ 处可微 $ \Leftrightarrow f(x) $ 在 $ x=x_0 $ 处可导，且
    $ \mathrm{d}x|_{x=x_0} =A(x_0)\Delta x=f'(x_0)\mathrm{d}x $ 
\end{Theo}

一般地，$ y=f(x) $ 则有 $ \mathrm{d}y=f'(x)\mathrm{d}x, $ 故导数 $ f'(x)=\frac{\mathrm{d}y}{\mathrm{d}x} $ 也称微商，即微分之商。

\subsection{微分公式与法则}

\begin{itemize}
    \item $ \mathrm{d}[f(x)\pm g(x)]=\mathrm{d}f(x)\pm\mathrm{d}g(x) $ ;
    \item $ \mathrm{d}[f(x)\cdot g(x)]=g(x)\mathrm{d}f(x)+f(x)\mathrm{d}g(x) $ ;
    \item $ g(x)\neq0 $ 时，$ \mathrm{d}\frac{f(x)}{g(x)}=\frac{g(x)\mathrm{d}f(x)-f(x)\mathrm{d}g(x)}{g^2(x)} $ .
\end{itemize}

\section{微分中值定理}

\subsection{罗尔定理}

\begin{Theo}[罗尔定理]

    若函数满足
    \begin{itemize}
        \item 在闭区间 $ [a,b] $ 连续；
        \item 在开区间 $ (a,b) $ 可导；
        \item $ f(a)=f(b), $ 
    \end{itemize}
    则至少存在一点 $ \xi\in(a,b) $ 使得 $ f'(\xi)=0. $ 
\end{Theo}
注意：
\begin{itemize}
    \item 常用于证明\textbf{形如}$ f(x)=g(x) $ 的方程在一开区间上有根；
    \item 难点：\begin{itemize}
        \item 构造区间；
        \item \textbf{构造辅助函数}，其中 $ F(x) $ 一般为 $ f(x)-g(x) $ 的原函数，或 $ f(x)-g(x) $ 的解。
    \end{itemize}
\end{itemize}

\begin{Field}[构造辅助函数的方法]

    \begin{enumerate}
        \item 将等式化为方程；
        \item 变形，直到能够积分为止；具体而言，可以\begin{itemize}
            \item 同乘同除一因式；
            \item 添加能正负相消的两项。
        \end{itemize}
        \item 取不定积分；
        \item 整理，直到满足其在 $ [a,b] $ 上连续，$ (a,b) $ 上可导为止。
    \end{enumerate}
    注意，以上方法仅用于发掘可能的辅助函数，不注重严谨性。
\end{Field}

不难发现，\begin{itemize}
    \item $ xf'(x) + nf(x) = 0 \Rightarrow F(x) = x^n f(x); $ 
    \item $ f(x)g'(x)-f'(x)=0\Rightarrow F(x)=f(x)e^{-g(x)} $ ;
    \begin{itemize}
        \item $ f(x)+f'(x)=0\Rightarrow F(x)=f(x)e^x $ ;
        \item $ f(x)-f'(x)=0\Rightarrow F(x)=f(x)e^{-x} $ ;
        \item $ f(x)+\lambda f'(x)=0\Rightarrow F(x)=f(x)e^{\lambda x} $ ;
    \end{itemize}

\end{itemize}

\subsection{拉格朗日中值定理}

\begin{Theo}[拉格朗日中值定理]

    若函数 $ f(x) $ 满足
    \begin{enumerate}
        \item 在 $ [a,b] $ 上连续；
        \item 在 $ (a,b) $ 上可导， 
    \end{enumerate}
    则至少存在一点 $ \xi \in (a,b) $ 使得 $ f'(\xi)=\dfrac{f(b)-f(a)}{b-a} 
    \Rightarrow f'(\xi)(b-a)=f(b)-f(a).$
\end{Theo}

注意：\begin{enumerate}
    \item 证明：令 $\dis F(x)=f(x)-\frac{f(b)-f(a)}{b-a}x $ 并且利用罗尔定理；
    \item 可证方程在开区间 $ (a,b) $ 上有根；
    \item 可证不等式；
    \item 出现一函数增量 $ \Rightarrow $ 利用拉格朗日中值定理；
    \item 联系 $ f(x) $ 与 $ f'(x) $ 的桥梁；
    \item 求极限时，若出现 $ f(b) - f(a), $ 可能应用拉格朗日得
    $ f'(\xi)(b-a), $ 再利用夹逼定理处理 $ \xi; $ 
    \item 与割线有关。
\end{enumerate}

\begin{Infer}[拉格朗日中值定理推论]

    \begin{itemize}
        \item 若 $ f(x) $ 在 $ (a,b) $ 内可导，且 $ f'(x)\equiv0, $ 则 $ f(x) $ 在 $ (a,b) $ 内为常数；
        \item 若 $ f(x) $ 在 $ (a,b) $ 内可导，且 $ f'(x)\equiv k, $ 则在 $ (a,b) $ 上 $ f(x)=kx+c, $ 
        其中 $ c $ 为常数。
    \end{itemize}
\end{Infer}

\subsection{柯西中值定理}

\begin{Theo}[柯西中值定理]

    设 $ f(x) $ 和 $ g(x) $ 满足
    \begin{enumerate}
        \item 在闭区间 $ [a,b] $ 上皆连续；
        \item 在开区间 $ (a,b) $ 内皆可导；
        \item $ g'(x)\neq 0, $ 
    \end{enumerate}
    则存在 $ \xi\in (a,b) $ 使得 $ \frac{f(b)-f(a)}{g(b)-g(a)}=\frac{f'(\xi)}{g'(\xi)}(a<\xi<b). $ 
\end{Theo}

注意：\begin{enumerate}
    \item 证明：令 $ F(x)=\frac{f(b)-f(a)}{g(b)-g(a)}g(x)-f(x), $ 利用罗尔定理可证；
    \item 可证方程有根；
    \item 出现两个函数增量 $ \Rightarrow $ 使用柯西中值定理证明；
    \item 若要求两个不同的中值，则寻求一特殊点 $ c $ 将 $ a,b $ 分为两子区间，并分别运用中值定理，如罗尔定理。
\end{enumerate}

\subsection{泰勒中值定理}

\begin{Theo}[带拉格朗日余项的泰勒中值定理]

    设 $ f(x) $ 在包含$ x_0 $ 的一邻域$ U_\delta(x_0) $ 内有 $ n+1 $ 阶导数，则对 $ x\in U_\delta, $ 
    有$$
        f(x) = \sum_{i=0}^n \frac{f^{(i)}(x_0)}{n!}(x-x_0)^n + R_n(x)
    $$ 
    其中，$\dis R_n(x)=\frac{f^{(n+1)}(\xi)}{(n+1)!}(x-x_0)^{n+1}, \xi $ 介于 $ x $ 与 $ x_0, $ 
    $ R_n(x) $ 称为拉格朗日余项。
\end{Theo}

\begin{Theo}[带皮亚诺余项的泰勒定理]

    设 $ f(x) $ 在 $ x=x_0 $ 处有 $ n $ 阶导数，则$$
        f(x) = \sum_{i=0}^n \frac{f^{(i)}(x_0)}{n!}(x-x_0)^n + R_n(x)
    $$ 其中 $ R_n(x)=o[(x-x_0)^n],\ x\rightarrow x_0 $ 称为皮亚诺余项。
\end{Theo}

以下是一些常用函数的带皮亚诺余项的马克劳林展开式。

\begin{itemize}
    \item $\dis e^x = \sum_{i=0}^n \dfrac{x^n}{n!}+ o(x^n)$;
    \item $\dis \cos x = 1 - \frac{x^2}{2} + \dfrac{x^4}{24} +\dots + \frac{(-1)^{n}x^{2n}}{(2n)!} + o(x^{2n})$ ;
    \item $\dis \sin x = x - \frac{x^3}{6} + \dots + \frac{(-1)^{n}x^{2n+1}}{(2n+1)!} + o(x^{2n+1}) $ ;
    \item $\dis \arcsin x = x + \frac{x^3}{6} +o(x^3)$ ;
    \item $\dis \tan x = x + \frac{x^3}{3} + o(x^3)$ ;
    \item $\dis \arctan x = x - \frac{x^3}{3} + o(x^3)$ ;
    \item $\dis \ln (1+x) = x - \frac{x^2}{2} + \frac{x^3}{3} + \dots + \frac{(-1)^{n-1}x^n}{n} + o(x^n) $;
    \item $\dis \ln(1-x) = -(x+\frac{x^2}{2} + \frac{x^3}{3}) + o(x^3)$;
    \item $\dis (1+x)^\alpha = 1+\sum_{k=1}^n C_\alpha^kx^k + o(x^n) $ ,
    其中 $\dis C_\alpha^k=\frac{\prod_{i = 0}^{k-1}(\alpha - i)}{k!} $ 

    如，$ \dis \sqrt{1+x} = 1 + \dfrac{1}{2}x - \dfrac{1}{8}x^2 + o(x^2); $ 
    \item $\dis \frac{1}{1-x} = \sum_{i=0}^n x^i + o(x^n) $ ;
    \item $\dis \frac{1}{1+x} = \sum_{i=0}^n (-1)^i x^i + o(x^n) $;
\end{itemize}

注意到泰勒展开时，对展开式的第二项，有
\begin{equation*}
    \begin{aligned}
        \sin x \xLongrightarrow{\textrm{正负变换}}
        \arcsin x \xLongrightarrow{\textrm{去掉阶乘}}
        \tan x \xLongrightarrow{\textrm{正负变换}}
        \arctan x \xLongrightarrow{\textrm{增加阶乘}}
        \sin x .
    \end{aligned}
\end{equation*}

注意，对于泰勒公式，
\begin{enumerate}
    \item \begin{enumerate}
        \item 可证不等式；
        \item 可证方程有根；
        \item 可求等价函数；
        \item 可求极限；
        \item 可求高阶导数值；
        \item 可近似计算；
    \end{enumerate}
    \item 在不等式证明、方程有根题目中应用拉格朗日余项，否则应用皮亚诺余项；
    \item 题目出现高阶(至少2阶)导数 $ \Rightarrow $ 可能应用泰勒中值定理；
    \item 一般在导数已知的点展开；
    \item 使用泰勒公式时，$ x\rightarrow 0 $ 可以推广为 $ \square \rightarrow 0. $ 
\end{enumerate}

\section{导数应用}

\subsection{函数单调判定定理}

\begin{Theo}[函数单调判定定理]

    设函数 $ f(x) $ 在 $ (a,b) $ 可导，若恒有 $ f'(x)>0(<0) $，则 $ f(x) $ 在 $ (a,b) $ 内单调递增(递减)；
    若恒有 $ f'(x)\geq0(\leq 0) $ ，则其在 $ (a,b) $ 内单调不减(不增)。
\end{Theo}

注意，\begin{itemize}
    \item 此为充分不必要条件；
    \item 成立条件：要求一区间，而非一点。反例：$$
        f(x)=\begin{cases}
            \frac{x}{2}+x^2\sin \frac{1}{x},& x\neq 0\\
            0,& x=0
        \end{cases}
    $$ 
    \item 可证不等式。
\end{itemize}

\section{极值}

\subsection{定义}

\begin{Def}[极值]

    设函数 $ f(x) $ 在 $ (a,b) $ 内有定义，$ x_0\in (a,b) $ ，则
    \begin{itemize}
        \item 若 $ x_0 $ 存在一邻域 $ U $ ，使得对 $ U $ 内任一点 $ x \neq x_0 $，
        都有 $ f(x)<f(x_0) $ ，则称 $ f(x_0) $ 为函数 $ f(x) $ 的一个极大值，
        $ x_0 $ 为 $ f(x) $ 一个极大值点；
        \item 若 $ x_0 $ 存在一邻域 $ U $ ，使得对 $ U $ 内任一点 $ x \neq x_0 $，
        都有 $ f(x)>f(x_0) $ ，则称 $ f(x_0) $ 为函数 $ f(x) $ 的一个极小值，
        $ x_0 $ 为 $ f(x) $ 一个极小值点；
    \end{itemize}
    函数的极大值和极小值统称极值。极大值点与极小值点统称极值点。
\end{Def}

\subsection{极值的性质}

\begin{itemize}
    \item 若 $ f'(x_0)=0 $ ，则称 $ x_0 $ 为一个驻点；
    \item 若 $ f(x) $ 在 $ x=x_0 $ 处可导，而$ f(x_0) $ 为极值，则 $ f'(x_0)=0, $
    即可导的极值点导数为零；
    \item 极值点 $ \nLeftrightarrow $ 驻点，如 $ f(x)=|x|, x=0; $ 
    \item 可能的极值点：\begin{itemize}
        \item $ f'(x)=0 $ 处；
        \item $ f'(x) $ 不存在处。
    \end{itemize}
\end{itemize}

\subsection{极值的充分条件}

\begin{Theo}[极值第一充分条件]

    设 $ f(x) $ 在 $ x_0 $ 处连续，在 $ 0<|x-x_0|<\delta $ 内可导，$ f'(x_0) $ 不存在或等于0，则
    \begin{enumerate}
        \item 若在 $ (x_0-\delta,x_0) $ 内任一点$ x $ 处有 $ f'(x)>0 $ ，而在 $ (x_0,x_0+\delta) $ 中
        任一点$ x $ 处有 $ f'(x)<0 $ ，则 $ f(x_0) $ 为极大值，$ x_0 $ 为极大值点；
        \item 若在 $ (x_0-\delta,x_0) $ 内任一点$ x $ 处有 $ f'(x)<0 $ ，而在 $ (x_0,x_0+\delta) $ 中
        任一点$ x $ 处有 $ f'(x)>0 $ ，则 $ f(x_0) $ 为极小值，$ x_0 $ 为极小值点；
        \item 若 $ (x_0-\delta,x_0) $ 内任一点与 $ (x_0,x_0+\delta) $ 中任一点
        处导数同号，则 $ f(x_0) $ 不为极值，$ x_0 $ 不为极值点。
    \end{enumerate}
\end{Theo}

\begin{Theo}[极值第二充分条件]

    设函数 $ f(x) $ 在 $ x_0 $ 处存在二阶导数，且 $ f'(x_0)=0 $ ， $ f^\pprime(x_0)\neq 0 $ ，
    则\begin{itemize}
        \item 当 $ f^\pprime(x_0)<0 $ 时，$ f(x_0) $ 为极大值，$ x_0 $ 为极大值点；
        \item 当 $ f^\pprime(x_0)>0 $ 时，$ f(x_0) $ 为极小值，$ x_0 $ 为极小值点。
    \end{itemize}
\end{Theo}

\section{最大值、最小值问题}

极值与最值的区别有\begin{enumerate}
    \item 前者局部，后者整体；
    \item 最值可为区间端点，极值不可以；
    \item 区间内部最值必为极值。
\end{enumerate}

最值可用于证明不等式。

应用问题的最值求法的步骤如下。
\begin{enumerate}
    \item 设出合适的自变量；
    \item 表示出目标函数；
    \item 求目标函数的最值。
\end{enumerate}

\subsection{凹向与拐点}

\begin{Def}[凹凸性]

    设 $ f(x) $ 在区间 $ I $ 连续，若对任意两点 $ x_1 \neq x_2 $ ，都有$$
        f(\frac{x_1+x_2}{2})>(<)\frac{1}{2}[f(x_1)+f(x_2)],
    $$ 则称 $ f(x) $ 在 $ I $ 上是凸(凹)的。
\end{Def}

几何上，若曲线 $ y=f(x) $ 上任意两点割线在曲线下(上)，则 $ y=f(x) $ 是凸(凹)的。

若 $ f(x) $ 有切线，而每一点的切线都在曲线之上(下)，则称 $ y=f(x) $ 是凸(凹)的。

\begin{Theo}[凹凸判定定理]

    设函数在 $ (a,b) $ 上有二阶导数，\begin{itemize}
        \item 若对任意 $ x\in (a,b) $ 都有 $ f^\pprime(x)>0 $ ，则称 $ f(x) $ 是凹函数；
        \item 若对任意 $ x\in (a,b) $ 都有 $ f^\pprime(x)<0 $ ，则称 $ f(x) $ 是凸函数；
    \end{itemize}
\end{Theo}

\begin{Def}[拐点]

    连续曲线上凹凸的分界点是曲线的拐点。
\end{Def}

\begin{Theo}[二阶可导点为拐点的必要条件]

    若 $ (x_0,f(x_0)) $ 为拐点，且 $ f^\pprime(x_0) $ 存在，则 $ f^\pprime(x_0)=0 $ 。
\end{Theo}

\begin{Theo}[二阶可导点为拐点的充分条件]

    设 $ y=f(x) $ 在 $ x=x_0 $ 连续，在 $ x_0 $ 的一去心邻域内二阶可导，\begin{enumerate}
        \item 若二阶导函数在 $ x=x_0 $ 两侧异号，则 $ (x_0,f(x_0)) $ 为一拐点；若同号，则不是拐点；
        \item 若在 $ x_0 $ 处二阶导数为零，而三阶导数不为零，则其为拐点；若三阶导数也为零，则不是拐点。
    \end{enumerate}
\end{Theo}

注意，$ f'(x) $ 的极值点事实上就是 $ f(x) $ 的拐点横坐标。

求拐点的步骤：\begin{enumerate}
    \item 求 $ x_0 $ 使得 $ f^\pprime(x_0)=0 $ 或不存在；
    \item 判断；
    \item 写拐点。
\end{enumerate}

\subsection{渐近线定义及其求法}

\sssubsection{水平渐近线}

若 $ {\displaystyle\lim_{x\rightarrow +\infty}}f(x)=b $
或 $ {\displaystyle\lim_{x\rightarrow -\infty}}f(x) = b, $ 
则称 $ y=b $ 是 $ f(x) $ 的一条水平渐近线。

\sssubsection{垂直渐近线}

若 $ {\displaystyle\lim_{x\rightarrow x_0^+}}f(x)=\infty $ 
或 $ {\displaystyle\lim_{x\rightarrow x_0^-}}f(x)=\infty, $ 则称 $ x=x_0 $ 是 $ f(x) $ 的一条垂直渐近线，也叫铅直渐近线。

垂直渐近线只需要讨论分母为零的点或者函数无定义的端点。

\sssubsection{斜渐近线} 

若 $ {\displaystyle\lim_{x\rightarrow +\infty}}f(x)-(kx+b)=0 $ $ (x\rightarrow-\infty), $ 则称
$ y=kx+b $ 为斜渐近线。

具体而言，若\begin{enumerate}
    \item $ {\displaystyle\lim_{x\rightarrow +\infty}}\frac{f(x)}{x}=k; $ 
    \item $ {\displaystyle\lim_{x\rightarrow +\infty}}f(x)-kx = b, $ 
\end{enumerate}
则有斜渐近线 $ y=kx+b $ 。$ x\rightarrow-\infty $ 时同理。

注意，
\begin{itemize}
    \item 即使 $ {\displaystyle\lim_{x\rightarrow \infty}} \frac{y}{x} $ 存在，斜渐近线也不一定存在；
    \item 一侧不会同时存在水平渐近线和斜渐近线。
\end{itemize}

\section{曲率及曲率半径}

曲率为$$
   k=\frac{|f^\pprime(x)|}{(1+(f'(x))^2)^{\frac{3}{2}}}
$$ 

曲率半径$ \rho $ 为 $ k $ 的倒数。
\chapter{一元函数积分学}

\section{不定积分}

\subsection{原函数}

\begin{Def}[原函数]

    设函数 $ f(x) $ 和 $ F(x) $ 在区间 $ I $ 上有定义，若$ F'(x) = f(x) $ 在区间 $ I $ 上成立，则
    称 $ F(x) $ 为 $ f(x) $在$ I $ 上的一个原函数。
\end{Def}

注意，原函数存在 $ \Rightarrow $ 有无穷多个，每个原函数之间仅差一个常数。

设 $ f(x) $ 在区间 $ I $ 上连续，则$ f(x) $ 在区间 $ I $ 上一定存在原函数，但是其不一定是
初等函数，如 $ \int \sin x^2 \mathrm{d}x $ ,$ \int \cos x^2 \mathrm{d}x $ ,
$ \int \frac{\sin x}{x} \mathrm{d}x $ ,$ \int \frac{\cos x}{x} \mathrm{d}x $ ,
$ \int \frac{1}{\ln x} \mathrm{d}x $ , $ \int e^{-x^2} \mathrm{d}x $ 等，称其为“积不出来”。

若$ f(x) $ 在区间 $ I $ 上存在第一类间断点或第二类中的无穷间断点，则其在该区间内没有原函数。

注意，函数在区间内有震荡间断点的，可能有原函数；总结如下。

\begin{Theo}[原函数存在定理]

    \begin{itemize}
        \item 连续函数存在原函数；
        \item 含有第一类间断点或无穷间断点的函数不存在原函数；
        \item 含有震荡间断点的函数可能存在原函数。
    \end{itemize}    
\end{Theo}

\subsection{不定积分}

\begin{Def}[不定积分]

    设 $ F(x) $ 为 $ f(x) $ 在区间 $ I $ 上的一个原函数，则$ f(x) $ 
    的所有原函数 $ \{F(x)+C\},C\in \mathbb{R} $ 为 $ f(x) $在$ I $ 的不定积分，
    记为 $\dis \int f(x) \mathrm{d}x$。
\end{Def}

\begin{Field}[不定积分基本性质]

    设 $\dis \int f(x)\mathrm{d}x = F(x)+C, $ 其中 $ F(x) $ 为 $ f(x) $ 一原函数，$ C $ 为一常数，
    则\begin{itemize}
        \item $\dis \int F'(x)\mathrm{d}x=F(x)+C $ 或 $\dis \int \mathrm{d}F(x)=f(x)+C ; $
        \item $\dis \left[\int f(x)\mathrm{d}x\right]'=f(x) $ 
        或 $\dis \mathrm{d}\left[\int f(x)\mathrm{d}x\right]=f(x)\mathrm{d}x ; $
        \item $\dis \int kf(x)\mathrm{d}x = k\int f(x)\mathrm{d}x ; $
        \item $\dis \int [f(x)\pm g(x)]\mathrm{d}x = \int f(x)\mathrm{d}x \pm \int g(x)\mathrm{d}x ; $
    \end{itemize}
\end{Field}

\sssubsection{基本积分公式}

$ C $ 为任意常数，则
\begin{enumerate}
    \item 幂函数
    
    $\dis \int x^\alpha \mathrm{d}x = \frac{x^{\alpha +1}}{\alpha + 1} + C\ (\alpha \neq -1), $
    
    特别地， $\dis \int \frac{1}{x} \mathrm{d}x = \ln|x| + C; $ 
    \item 指数函数
    
    $\dis \int a^x \mathrm{d}x = \frac{a^x}{\ln a} + C\ (a>0,a\neq 1) ; $
    
    特别地，$\dis \int e^x\mathrm{d}x = e^x+C ; $
    \item 对数函数
    
    $ \dis \int\ln x\mathrm{d}x = x \ln x - x + C; $ 
    \item 三角函数
    
    $ \dis \int \cos x \mathrm{d}x = \sin x + C, \dis \int \sin x \mathrm{d}x = -\cos x + C ; $

    $ \dis \int \tan x\mathrm{d}x = -\ln|\cos x| + C,\int \cot x \mathrm{d}x = \ln|\sin x| + C; $ 

    $ \dis \int \sec x\mathrm{d}x = \ln|\sec + \tan x| + C,\int \csc x= -\ln|\csc x + \cot x| + C
    = \ln|\csc x - \tan x| + C; $ 

    $\dis \int \frac{1}{\cos^2 x}\mathrm{d}x = \tan x + C, \dis \int -\frac{1}{\sin^2 x}\mathrm{d}x = \cot x + C ; $

    $\dis \int \frac{\tan x}{\cos x} \mathrm{d}x = \frac{1}{\cos x} + C, 
    \dis \int \frac{\cot x}{\sin x}\mathrm{d}x = -\frac{1}{\sin x} + C ; $
    \item 平方和
    
    $\dis \int \frac{1}{1+x^2}\mathrm{d}x = \arctan x + C = -\textrm{arccot} x + C; $

    $ \dis \int \dfrac{\mathrm{d}x}{a^2 + x^2} = \dfrac{1}{a}\arctan \dfrac{x}{a} + C; $    
    \item 平方差
    
    $ \dis \int \dfrac{\mathrm{d}x}{1-x^2} = \dfrac{1}{2}\ln\left|\dfrac{1+x}{1-x}\right| + C; $ 

    $ \dis \int \dfrac{\mathrm{d}x}{a^2-x^2} = \dfrac{1}{2a}\ln\left|\dfrac{a+x}{a-x}\right| + C; $ 
    
    \item 根号下平方差
    
    $ \dis \int \dfrac{\mathrm{d}x}{\sqrt{1-x^2}} = \arcsin x + C,$

    $ \dis \int \dfrac{\mathrm{d}x}{\sqrt{a^2-x^2}} = \arcsin \dfrac{x}{a} + C,$
    \item 根号下平方和
    
    $ \dis \int \dfrac{\mathrm{d}x}{\sqrt{a^2+x^2}} = \ln(x + \sqrt{a^2+x^2}) + C,$

    $ \dis \int \dfrac{\mathrm{d}x}{\sqrt{x^2 - a^2}} = \ln\left|x - \sqrt{x^2-a^2}\right| + C;$
\end{enumerate}    



\subsection{不定积分法}

\begin{Field}[第一换元法(凑微分)]

    设 $ g(\varphi(x)) = f(x) $ ，则计算的原则是
    \begin{equation*}
        \begin{aligned}
            \int f(x)\mathrm{d}x &= \int g(\varphi(x))\varphi'(x)\mathrm{d}x \\
            &= \int g(\varphi(x))\mathrm{d}\varphi(x) \\ 
            &= F(\varphi(x))+ C
        \end{aligned}
    \end{equation*}
\end{Field}

常见的凑微分形式如下。

\begin{enumerate}
    \item $ \int f(ax+b)(ax+b)\mathrm{d}x = \frac{1}{a}\int f(ax+b)\mathrm{d}(ax+b) (a\neq 0) ; $
    \item $ \int f(ax^n+b)x^{n-1}\mathrm{d}x = \frac{1}{na}\int f(ax^n+b)\mathrm{d}(ax^n + b) (an\neq0) ; $
    \item $ \int \frac{f(\ln x)}{x}\mathrm{d}x = \int f(\ln x)\mathrm{d}(\ln x); $
    \item $ \int f(\frac{1}{x})\frac{1}{x^2}\mathrm{d}x = -\int f(\frac{1}{x})\mathrm{d}(\frac{1}{x}) ; $
    \item $ \int \frac{f(\sqrt x)}{\sqrt x} \mathrm{d}x = 2 \int f(\sqrt x)\mathrm{d}(\sqrt x) ; $
    \item $ \int f(a^x)a^x\mathrm{d}x = \frac{1}{\ln a}\int f(a^x)\mathrm{d}(a^x) $,
    特别地，$ \int f(e^x)e^x\mathrm{d}x = \int f(e^x)\mathrm{d}(e^x) ; $
    \item $ \int f(\sin x)\cos x \mathrm{d}x = \int f(\sin x)\mathrm{d}(\sin x) ; $
    \item $ \int f(\cos x)\sin x \mathrm{d}x = -\int f(\cos x)\mathrm{d}(\cos x); $
    \item $ \int \frac{f(\tan x)}{\cos^2 x}\mathrm{d}x = \int f(\tan x)\mathrm{d}(\tan x) ; $
    \item $ \int \frac{f(\cot x)}{\sin^2 x}\mathrm{d}x = -\int f(\cot x)\mathrm{d}(\cot x) ; $
    \item $ \int f(\sec x)\sec x \tan x \mathrm{d}x = \int f(\sec x)\mathrm{d}(\sec x) ; $
    \item $ \int f(\csc x)\csc x \cot x\mathrm{d}x = -\int f(\csc x)\mathrm{d}(\csc x) ; $
    \item $ \int \frac{f(\arcsin x)}{\sqrt{1-x^2}}\mathrm{d}x = \int f(\arcsin x)\mathrm{d}(\arcsin x) ; $
    \item $ \int \frac{f(\arccos x)}{\sqrt{1-x^2}}\mathrm{d}x = -\int f(\arccos x)\mathrm{d}(\arccos x) ; $
    \item $ \int \frac{f(\arctan x)}{1+x^2}\mathrm{d}x = \int f(\arctan x)\mathrm{d}(\arctan x) ; $
    \item $ \int \frac{f(\textrm{arccot} x)}{1+x^2}\mathrm{d}x = -\int f(\textrm{arccot} x)\mathrm{d}(\textrm{arccot} x) ; $
    \item $ \int \frac{f(\arctan \frac{1}{x})}{1+x^2}\mathrm{d}x = -\int f(\arctan \frac{1}{x})\mathrm{d}(\arctan \frac{1}{x}) ; $
    \item $ \int \frac{f|ln(x+\sqrt{x^2+a^2})|}{\sqrt{x^2+a^2}}\mathrm{d}x = 
    \int f|ln(x+\sqrt{x^2+a^2})|\mathrm{d}(\ln(x+\sqrt(x^2+a^2))) ; $
    \item $ \int \frac{f|ln(x+\sqrt{x^2-a^2})|}{\sqrt{x^2-a^2}}\mathrm{d}x = 
    \int f|ln(x+\sqrt{x^2-a^2})|\mathrm{d}(\ln(x+\sqrt(x^2-a^2))) ; $
    \item $ \int \frac{f'(x)}{f(x)}\mathrm{d}x = \ln|f(x)|+C\ (f(x)\neq0) $ .
\end{enumerate}

\begin{Field}[第二换元法]

    设 $ x=\varphi(t) $ 具有连续导数且单调，且 $ \varphi'(t)\neq 0 $ ，若
    $$
        \int f[\varphi(t)]\varphi'(t)\mathrm{d}t=G(t)+C
    $$ 
    则有$$
        \int f(x)\mathrm{d}x 
        \xlongequal{\textrm{令}x=\varphi(t)}
        \int f[\varphi(t)]\varphi'(t)\mathrm{d}t = G(t) + C = G[\varphi^{-1}(x)]+C
    $$ 
    其中 $ t=\varphi^{-1}(x) $ 为 $ x=\varphi(t) $ 的反函数。
\end{Field}

\begin{comment}
被积函数是 $ x $ 与 $ \sqrt[n]{ax+b} $ 或 $ x $ 与 $ \dis\sqrt[n]{\frac{ax+b}{cx+d}} $ 
    或由 $ e^x $ 构成的代数式的根式，则令 $ \sqrt[n]{g(x)}=t $，使得 $ x=\varphi(t) $ 中不再含有根式，
    此时做变量代换 $ x = \varphi(t) $ 即可；
    被积函数含有 $ \dis\sqrt{Ax^2+Bx+C}(A\not=0) $ 时，先根据$ A $ 的符号将其整理为
    $ \sqrt{A[(x-x_0)^2\pm l^2]} $ 或 $ \sqrt{-A[l^2-(x-x_0)^2]} $ ，然后做三角替换。   
\end{comment}

第二换元法多用于根式的被积函数，通过换元法去根式，具体而言，可以分为两类。
\begin{enumerate}
    \item 根号下一次方
    
    \begin{itemize}
        \item $ \dsqrt[n]{ax+b} = t; $ 
        \item $ \dsqrt[n]{\dfrac{ax+b}{cx+d}} = t; $ 
        \item 当式中同时存在 $ \dsqrt[n]{ax+b},\dsqrt[m]{ax+b} $ 时，取 $ t = \dsqrt[l]{ax+b}, $ 
        其中 $ l $ 为 $ m,n $ 的最小公倍数。
    \end{itemize}
    \item 根号下二次方
    
    \begin{itemize}
        \item 根式形如 $ \sqrt{a^2-x^2} $ 时，令 $ x=a\sin t, $ 此时 $ t\in[-\frac{\pi}{2},\frac{\pi}{2}]; $ 
        \item 根式形如 $ \sqrt{x^2-a^2} $ 时，令 $ x=a\sec t, $ 此时 $ t\in[0,\frac{\pi}{2})\cup(\frac{\pi}{2},\pi]; $ 
        \item 根式形如 $ \sqrt{a^2+x^2} $ 时，令 $ x=a\tan t, $ 此时 $ t\in (-\frac{\pi}{2},\frac{\pi}{2}). $ 
    \end{itemize}
\end{enumerate}

\begin{Field}[分部积分法]

    对 $ \int uv'\mathrm{d}x $ ，有$$
        \int uv'\mathrm{d}x = \int u\mathrm{d}v=uv-\int v\mathrm{d}u
    $$ 
\end{Field}

注意，\begin{itemize}
    \item 被积函数为两类函数乘积；
    \item 凑微分的优先性：指数 $ > $ 三角函数 $ > $ 幂函数。
\end{itemize}

注意，不定积分结果中一定包含常数 $ C, $ \textbf{即使在计算过程中积分号下内容相消干净。}

分部积分法可以利用列表法做。具体而言，为上导下积，对角相乘，正负相间。
以 $ \dis \int x^2e^x \mathrm{d}x $ 为例，则有表格

\begin{center}
    \begin{tabular}{c|ccc}
        $ x^2 $ & $ 2x $ & $ 2 $ & $ 0 $ \\\hline 
        $ e^x $ &$ e^x $ &$ e^x $ &$ e^x $ 
    \end{tabular}    
\end{center}
    
此时结果为  $ x^2 e^x - 2x e^x + 2 e^x + C. $ 

\subsection{特殊类型的不定积分法}

\subsubsection{通用方法}

\begin{itemize}
    \item 倒代换 - 分母次数高的，令$ x = \dfrac{1}{t}; $ 
    \item 整体代换 - 含有复杂函数的，令复杂部分为 $ t; $
\end{itemize}

\subsubsection{有理分式的不定积分}

\begin{Theo}[]

    任何实系数多项式在实数域内均可分解为一次因式和二次因式的乘积。
\end{Theo}

\begin{Theo}[]

    设 $ f(x)=\frac{P_m(x)}{Q_n(x)}$ 是有理真分式，若在实数域内分母 $ Q_n(x) $ 可因式分解为 $$
        Q_n(x)=(x-a)^\alpha(x^2+px+q)^\beta
    $$ 
    则有$$
        \frac{P_m(x)}{Q_n(x)}=\sum_{i=1}^\alpha \frac{A_i}{(x-a)^i}+
        \sum_{i=1}^\beta \frac{C_ix+D_i}{(x^2+px+q)^i}
    $$ 其中 $ A_i\ B_i\ C_i\ D_i $ 等系数都是实数。
\end{Theo}

具体而言，其步骤为
\begin{enumerate}
    \item 将假分式分为多项式和真分式；
    \item 将真分式分母整理至最简；
    \item 裂项，分别积分。
\end{enumerate}

求待定系数时，去分母，并注意\begin{itemize}
    \item 多项式相等条件：同次数的系数相等；
    \item 待定系数的个数为分母次数。
\end{itemize}

分母次数大于$ 4 $ 的时候，一般不用裂项法，而是寻求特殊方法。

待定系数时可以代入特殊值求解系数。

一个例子是，
\begin{equation*}
    \begin{aligned}
        \dfrac{x^4+2x^3+3x^2+4x+1}{(x-1)(x^2-1)(x^2+1)}=&\dfrac{x^4+2x^3+3x^2+4x+1}{(x+1)(x-1)^2(x^2+1)}\\ 
        =& \dfrac{A}{x+1} + \dfrac{B}{x-1} + \dfrac{C}{(x-1)^2} + \dfrac{Dx+E}{x^2+1}
    \end{aligned}
\end{equation*}

注意，
\begin{itemize}
    \item 分母有高次幂的，待定系数时次数从低到高都要有；
    \item 待定系数时分子次数比分母次数低一次。
\end{itemize}

\subsubsection{三角有理函数的不定积分}

三角有理函数指由 $ \sin x $ 和 $ \cos x $ 进行有限次四则运算得到的函数。

\begin{Field}[三角有理函数的不定积分一般方法]

    利用万能代换，即$\tan \dis\frac{x}{2} = t $，此时 $ \sin x = \dis\frac{2t}{t^2+1} $ ,
    $ \cos x = \dis\frac{t^2-1}{t^2+1} $ , $ x = 2\arctan t $ ,
    $ \dis\frac{\mathrm{d}x}{\mathrm{d}t} = \displaystyle\frac{2}{1+t^2} $ .
\end{Field}

注意，应用万能代换时， $ x $ 的次数不应当超过一次，形如$$
    \int \dis\frac{a_1\sin x + b_1 \cos x + c_1}{a_2\sin x + b_2 \cos x + c_2}
$$ 

\sssubsection{特殊三角有理函数的积分}

\begin{itemize}
    \item 对 $$
        I = \int \frac{a\sin x + b\cos x}{c\sin x+d\cos x}\mathrm{d}x
    $$ 
    可以使用待定系数法，令$$
        a\sin x+b\cos x = A(c\sin x+d\cos x)+B(c\sin x+ d\cos x)'
    $$ 
    其中$ A,B $ 是待定系数，将其解出并代入原式，得$$
        I=Ax+B\ln|c\sin x+d\cos x|+C
    $$ 
    \item \begin{itemize}
        \item 若$ R(-\sin x, \cos x)=-R(\sin x, \cos x) $ ，令 $ \cos x = t $ 并换元；
        \item 若$ R(\sin x, -\cos x)=-R(\sin x, \cos x) $ ，令 $ \sin x = t $ 并换元；
        \item 若$ R(-\sin x, -\cos x)=R(\sin x, \cos x) $ ，令 $ \tan x = t $ 并换元；
        \item 针对上一条，若$ R $ 形如$$
            \int \sin^{2m} x + \cos^{2n} x \mathrm{d}x
        $$ 其中 $ m,n\in N $ ，则应使用二倍角公式降幂。
    \end{itemize}
\end{itemize}


\section{定积分}

\subsection{定积分概念}

\begin{Def}[]

    将 $ [a,b] $ 任意地划分为 $ n $ 个小区间 $ a = x_1<x_2<\dots<x_{n-1}<x_n=b $ ，
    细度 $ d=\max_{1\leq i \leq n}(x_i-x_{i-1}) $ , $ \xi_i $ 为 $ [x_{i-1},x_i] $ 任意一点；
    则$ f(x) $在$ [a,b] $ 上的定积分为$$
        \int_a^b f(x)\mathrm{d}x = {\displaystyle\lim_{d\rightarrow 0}}
        \sum_{i=i}^nf(\xi_i)\Delta x_i
    $$ 
    若上述极限存在。
    若$ f(x) $在$ [a,b] $ 上确有定积分，则称 $ f(x) $在$ [a,b] $ 可积。
\end{Def}

\begin{Theo}[连续必可积]

    若$ f(x) $ 在闭区间 $ [a,b] $ 连续，则$ f(x) $在$ [a,b] $ 可积。
\end{Theo}

\begin{Theo}[]

    若$ f(x) $在$ [a,b] $ 有界，且在 $ [a,b] $ 中仅有有限个间断点，则$ f(x) $在$ [a,b] $ 可积。
\end{Theo}

\begin{Theo}[可积必有界]

    若$ f(x) $在$ [a,b] $ 可积，则其在 $ [a,b] $ 有界。
\end{Theo}

\begin{Theo}[定积分值的变量符号无关性]

    定积分的值与积分变量的符号无关，即$$
        \int_a^b f(x)\mathrm{d}x=\int_a^b f(t)\mathrm{d}t
    $$ 
\end{Theo}

\sssubsection{利用定积分定义求极限}

\begin{equation*}
    \begin{aligned}
        \int_a^b f(x)\mathrm{d}x&={\displaystyle\lim_{d\rightarrow 0}}\sum_{i=i}^nf(\xi_i)\Delta x_i\\
        &={\displaystyle\lim_{n\rightarrow \infty}}\sum_{i=1}^n f(a+\frac{i}{n}(b-a))\frac{b-a}{n}
    \end{aligned}
\end{equation*}

具体而言，其步骤为
\begin{itemize}
    \item 令 $ \Delta x_i $ 为 $ \frac{1}{n} $ 或 $ \frac{k}{n} ; $
    \item 取 $ \xi_i $ 的位置，如\begin{itemize}
        \item $ \xi_i $ 取右端点： $ a+\frac{i}{n}(b-a) ; $
        \item $ \xi_i $ 取左端点： $ a+\frac{i-1}{n}(b-a) ; $
        \item $ \xi_i $ 取中点： $ a+\frac{2i-1}{2n}(b-a) ; $
        \item $ \xi_i $ 取三等分点： $ a+\frac{3i-2}{3n}(b-a) ; $
        \item $ \xi_i $ 取三分之二等分点： $ a+\frac{3i-1}{3n}(b-a) $ 等等。
    \end{itemize}
    \item 求 $ a={\displaystyle\lim_{n\rightarrow \infty}}\xi_1,b={\displaystyle\lim_{n\rightarrow \infty}}\xi_n $
    \item 将求得的 $ a,b $ 用于验证 $ \Delta x_i $ 的取法是否正确，
    即是否有$ \sum_{i=1}^n \Delta x_i = b-a $ 。
\end{itemize}

特别地，有
$$
    {\displaystyle\lim_{n\rightarrow \infty}}\sum_{i=1}^n f(\dfrac{i}{n})\dfrac 1n =\int_0^1f(x)\mathrm{d}x
$$ 
此即所谓“零到一，$n $ 等分，取端点”的极限。

\sssubsection{利用二重积分定义求极限}

若 $ f(x,y) $ 在 $ [a,b]\times [c,d] $ 上连续，则有$$
    {\displaystyle\lim_{n\rightarrow \infty}}\sum_{i=1}^n\sum_{j=1}^n f(\xi_i,\eta_j)\dfrac{b-a}{n}
    \dfrac{d-c}{n} = \int_a^b\mathrm{d}x\int_c^df(x,y)\mathrm{d}y
$$ 
其中$$
    \xi_i\in [a+\dfrac{i-1}{n}(b-a),a+\dfrac{i}{n}(b-a)],
    \xi_j\in [c+\dfrac{j-1}{n}(d-c),a+\dfrac{j}{n}(d-c)],
$$ 
其具体的步骤与一维情况类似。

特别地，有
$$
    {\displaystyle\lim_{n\rightarrow \infty}}\sum_{i=1}^n\sum_{j=1}^n f(\dfrac{i}{n},\dfrac{j}{n})
    \dfrac{1}{n^2} = \int_0^1\mathrm{d}x\int_0^1f(x,y)\mathrm{d}y
$$ 


\sssubsection{$ n $ 项和求法}

$ n $ 项和的求和有如下的方式。
\begin{itemize}
    \item 直接求和；
    \item 夹逼定理；
    \item 定积分定义；
    \item 数项级数。
\end{itemize}

\begin{Field}[定积分的几何意义]

    定积分的几何意义为面积的代数和。
\end{Field}

\subsection{定积分的性质}

\begin{itemize}
    \item 线性性 - $\dis \int_a^b [k_1 f_1(x)]+k_2[f_2(x)]\mathrm{d}x = 
    k_1\int_a^b f_1(x)\mathrm{d}x + k_2\int_a^b f_2(x)\mathrm{d}x ; $
    \item 区间可加性 - $\dis \int_a^b f(x)\mathrm{d}x = \int_a^c f(x)\mathrm{d}x + \int_c^b f(x)\mathrm{d}x, $
    即使 $ c\not\in [a,b] ; $
    \item 比较定理 - 设 $ a\leq b, f(x)\leq g(x), $ 则 $\dis \int_a^b f(x)\mathrm{d}x \leq \int_a^b g(x)\mathrm{d}x ; $
    
    推论 - $\dis a<b \Rightarrow |\int_a^b f(x) \mathrm{d}x | < \int_a^b|f(x)|\mathrm{d}x ; $
    \item 估值定理 - $\dis a < b, m\leq f(x)\leq M \Rightarrow m(b-a)\leq\int_a^b f(x)\mathrm{d}x \leq M(b-a); $
    \item 积分中值定理 - 设 $ f(x) $在$ [a,b] $ 上连续，则$\exists \xi \in [a,b]  $ 
    使得 $\dis \int_a^b f(x)\mathrm{d}x = f(\xi)(b-a) ; $

    加强形式 - 设 $ f(x) $在$ [a,b] $ 上连续，则$\exists \xi \in (a,b)  $ 
    使得 $\dis \int_a^b f(x)\mathrm{d}x = f(\xi)(b-a) ; $

    加强形式 - 设 $ f(x),g(x) $在$ [a,b] $ 上连续，$ g(x) $ 不变号，则$\exists \xi \in (a,b)  $ 
    使得 $\dis \int_a^b f(x)g(x)\mathrm{d}x = f(\xi)\int_a^bg(x)\mathrm{d}x ; $
\end{itemize}

注意，\begin{itemize}
    \item 定积分值与被积函数在有限个点上的值无关；
    \item 定积分中值定理可证方程有根。
\end{itemize}

\subsection{重要定理、公式、关系}

\begin{Def}[变上限函数]

    设 $ f(x) $在$ [a,b] $ 可积，则函数形如 $ F(x)=\int_a^x f(t)\mathrm{d}t, x\in [a,b]  $ 
    称为变上限积分函数。
\end{Def}

\begin{Theo}[变上限积分求导定理]

    若$ f(x) $在$ [a,b] $ 可积，则$ F(x) $在$ [a,b] $ 连续。

    若$ f(x) $在$ [a,b] $ 连续，则$ F(x) $在$ [a,b] $ 上可导，且 $ F'(x) = f(x) $。
\end{Theo}

注意，\begin{itemize}
    \item 设 $$
        F(x)=\int_{\varphi_1(x)}^{\varphi_2(x)} f(t)\mathrm{d}t
    $$ 且 $ \varphi_1(x),\varphi_2(x) $ 可导，$ f(x) $ 连续，则有$$
        F'(x) = f(\varphi_2(x))\varphi_2'(x)-f(\varphi_1(x))\varphi_1'(x)
    $$ 

    变限积分求导，被积函数不含 $ x. $ 
    \item 若$ x_0 $ 为 $ f(x) $在$ [a,b] $ 上的一个跳跃间断点，则$ F(x) $在$ x_0 $ 连续，
    但不可导；
    \item 若$ x_0 $ 为 $ f(x) $在$ [a,b] $ 上的一个可去间断点，则$ F(x) $在$ x_0 $ 处可导，
    但 $ F(x) $ 不是$ f(x) $ 的原函数。
\end{itemize}

\begin{Theo}[牛顿——莱布尼茨公式]

    设$ f(x) $在$ [a,b] $可积，$ F(x) $为$ f(x) $在$ [a,b] $上的任意原函数，
    则有$$
        \int_a^b f(x)\mathrm{d}x=F(x)\Big|_a^b=F(b)-F(a)
    $$ 
\end{Theo}

\subsection{定积分求法}

\sssubsection{利用牛顿莱布尼茨公式}

$$
   \int_a^b f(x)\mathrm{d}x=F(x)\Big|_a^b=F(b)-F(a)
$$ 

\sssubsection{分部积分法}

$$
    \int_a^b u'v \mathrm{d}x = \int_a^b u\mathrm{d}v = uv\Big|_a^b - \int_a^b v\mathrm{d}u
$$ 

其利用技巧与不定积分的分部积分法完全一致。

\sssubsection{换元积分法}

设 $ f(x) $ 在 $ [a,b] $ 上连续，若变量替换 $ x=\varphi(t) $ 满足
\begin{enumerate}
    \item $ \varphi'(t) $ 在 $ [\alpha,\beta] $ (或者$ [\beta,\alpha] $ ) 上连续；
    \item $ \varphi(\alpha) = a , \varphi(\beta) = b $ ，且当 $ \alpha < t < \beta $ 时，
    有 $ a \leq \varphi(t) < b, $ 
\end{enumerate}
则有$$
    \int_a^b f(x)\mathrm{d}x = \int_\alpha^\beta f[\varphi(t)]\varphi'(t)\mathrm{d}t
$$ 

\subsection{常用的定积分公式}

\sssubsection{对称区间奇偶函数的积分公式}

\begin{enumerate}
    \item 设 $ f(x) $ 是在区间 $ [-a,a],(a>0) $ 上连续的偶函数，则$$
        \int_{-a}^a f(x)\mathrm{d}x = 2 \int_0^a f(x)\mathrm{d}x
    $$ 
    \item 设 $ f(x) $ 是在区间 $ [-a,a],(a>0) $ 上连续的奇函数，则$$
        \int_{-a}^af(x)\mathrm{d}x = 0
    $$ 
\end{enumerate}

可用于积分等式的方法：
\begin{enumerate}
    \item 含有中值：积分中值定理；
    \item 不含中值：\begin{itemize}
        \item 积分区间分割；
        \item 换元法
    \end{itemize}
\end{enumerate}

\sssubsection{周期函数的积分公式}

设 $ f(x) $ 在 $ (-\infty,+\infty) $ 内是以 $ T $ 为周期的周期函数，则对任意常数 $ a $ 、任意自然数 $ n $ 
都有\begin{enumerate}
    \item $ \dis \int_a^{a+T}f(x)\mathrm{d}x = \int_0^T f(x)\mathrm{d}x ; $
    \item $ \dis \int_a^{a+nT}f(x)\mathrm{d}x = n\int_0^T f(x)\mathrm{d}x ; $
\end{enumerate}

\sssubsection{三角函数的积分公式}

\begin{Theo}[Wallis公式 - 分部积分推导]

    $$
        \int_0^\frac{\pi}{2} \sin^{2n} x \mathrm{d} x = \int_0^\frac{\pi}{2} \cos^{2n} x \mathrm{d} x = 
        \dfrac{(2n-1)!!}{(2n)!!}\cdot \dfrac{\pi}{2}
    $$ 
    $$
        \int_0^\frac{\pi}{2} \sin^{2n+1} x \mathrm{d} x = \int_0^\frac{\pi}{2} \cos^{2n+1} x \mathrm{d} x = 
        \dfrac{(2n)!!}{(2n+1)!!}
    $$ 
    $$
        \int_0^\frac{\pi}{2} \sin^{n} x \mathrm{d} x = \int_0^\frac{\pi}{2} \cos^{n} x \mathrm{d} x = 
        \dfrac{(n-1)!!}{(n)!!}\cdot (\dfrac{\pi}{2})^\texttt{(int)(!(n\%2))}
    $$ 
\end{Theo}

注意，当被积函数中有自然数 $ n $ 时，通过分部积分法构造递推公式，并用递推公式推出结果。

三角函数的公式还有如下几个。设 $ f(x) $ 在 $ [0,1] $ 连续，则
\begin{itemize}
    \item $ \dis \int_0^\frac{\pi}{2} f(\sin x)\mathrm{d}x=\int_0^\frac{\pi}{2} f(\cos x)\mathrm{d}x $ 
    \item $\dis \int_0^\pi xf(\sin x)\mathrm{d}x = \dfrac{\pi}2\int_0^\pi f(\sin x)\mathrm{d}x = 
    \pi\int_0^\frac{\pi}{2} f(\sin x)\mathrm{d}x $ 
    \item $\dis \int_0^{\pi} f(\sin x)\mathrm{d}x = 2\int_0^\frac{\pi}{2} f(\sin x)\mathrm{d}x $ 
\end{itemize}

其中，$ ii) $ 利用了区间再现公式，具体而言，令 $ t = a + b - x $ 即可。

难以计算原函数的积分称为变态积分。对变态积分的计算，可以
\begin{itemize}
    \item 构造递推公式；
    \item 积分区间再现。
\end{itemize}

\section{反常积分}

定积分存在两种限制，由此可以引申出两种反常积分。
对 $ \dis \int_a^b f(x)\mathrm{d}x, $ 
\begin{itemize}
    \item $ [a,b] $ 必须是有限区间 $ \Rightarrow $ 无穷区间上的反常积分；
    \item $ f(x) $ 必须有界 $ \Rightarrow $ 无界函数的反常积分。
\end{itemize}

反常积分时定积分的极限。

\subsection{无穷区间上的反常积分}

$ \dis \int_a^{+\infty}f(x)\mathrm{d}x = {\displaystyle\lim_{t\rightarrow +\infty}}
\int_a^t f(x)\mathrm{d}x $ 称为无穷区间上的反常积分，
若其存在，则称该反常积分收敛，否则发散。\vspace{3pt}

对 $ \dis \int_{-\infty}^b f(x)\mathrm{d}x = {\displaystyle\lim_{t\rightarrow -\infty}}
\int_t^b f(x)\mathrm{d}x $ 同理。\vspace{3pt}

对 $ \dis \int_{-\infty}^{+\infty}f(x)\mathrm{d}x $，将其拆成以上两种情况的反常积分，若二者都收敛，
则称其收敛，否则称其发散。

注意，
\begin{enumerate}
    \item 计算方法与积分类似，如 $\dis \int_a^{+\infty} = F(x)\Big|^{+\infty}_a = 
    {\displaystyle\lim_{t\rightarrow +\infty}}F(t) - F(a); $ 
    \item 反常积分收敛时，适用对称奇偶性，否则不适用；
    \item $ p $ 积分：对 $ \dis \int_a^{+\infty}\dfrac{1}{x^p}\mathrm{d}x,\ a>0 $ ,
    当 $ p > 1 $ 时积分收敛，$ p\leq 1 $ 时发散；
    \item 利用分部积分法时，可以先进行不定积分，然后再引入上下限。
\end{enumerate}

\begin{Theo}[比较判别法]

    设 $ f $ 与 $ g $ 在 $ [a,+\infty) $ 上非负且在任意区间 $ [a,b],b<\infty $ 可积，且 $ f \leq g $ ，
    $ x\in[a,+\infty) $ ，则\begin{enumerate}
        \item 若 $ \int_a^{+\infty}g(x)\mathrm{d}x $ 收敛，则$ \int_a^{+\infty}f(x)\mathrm{d}x $收敛；
        \item 若 $ \int_a^{+\infty}f(x)\mathrm{d}x $ 发散，则$ \int_a^{+\infty}g(x)\mathrm{d}x $发散。
    \end{enumerate}
\end{Theo}

\begin{Theo}[比较判别法]

    设 $ f $ 与 $ g $ 在 $ [a,+\infty) $ 上非负，则
    \begin{enumerate}
        \item 若 $ g(x) > 0 $ 且 $ \dis {\displaystyle\lim_{x\rightarrow +\infty}}
        \dfrac{f(x)}{g(x)} = k \neq 0 $ ，则 $ \int_a^{+\infty}f(x)\mathrm{d}x $ 与
        $ \int_a^{+\infty}g(x)\mathrm{d}x $ 的敛散性相同；
        \item 若当 $ x\rightarrow+\infty $时，$ f\sim g $，则 $ \int_a^{+\infty}f(x)\mathrm{d}x $ 与
        $ \int_a^{+\infty}g(x)\mathrm{d}x $ 的敛散性相同。
    \end{enumerate}
\end{Theo}

\begin{Theo}[绝对收敛必收敛]

    若 $ f(x) \in C[a,\infty) $ 且 $ \dis \int_a^{+\infty}|f(x)|\mathrm{d}x $ 收敛，则
    $ \dis \int_a^{+\infty}f(x)\mathrm{d}x $ 收敛。
\end{Theo}

\subsection{无界函数的反常积分}

设 $ {\displaystyle\lim_{x\rightarrow b^-}}f(x) = \infty $，则称 $ b $ 为瑕点，
$ \dis \int_a^b f(x)\mathrm{d}x = {\displaystyle\lim_{t\rightarrow b^-}}
\int_a^t f(x)\mathrm{d}x $ 为无界函数的反常积分，若其存在，称其收敛，否则其发散。

瑕点在区间左端点或区间内时类似。瑕点在区间内而拆成两个积分时，二积分全部收敛则收敛，
否则发散。

注意，
\begin{enumerate}
    \item 计算方法与积分类似，如 $\dis \int_a^{b} = F(x)\Big|^{b}_a = 
    {\displaystyle\lim_{t\rightarrow b^-}}F(t) - F(a) $ ；
    \item 反常积分收敛时，适用对称奇偶性，否则不适用；
    \item $ p $ 积分：对 $ \dis \int_0^{a}\dfrac{1}{x^p}\mathrm{d}x,\ a>0 $ ,
    当 $ p < 1 $ 时积分收敛，$ p\geq 1 $ 时发散；
    \item 利用分部积分法时，可以先进行不定积分，然后再引入上下限。
\end{enumerate}

\begin{Theo}[比较判别法]

    设 $ {\displaystyle\lim_{x\rightarrow b^-}}f(x) = \infty,{\displaystyle\lim_{x\rightarrow b^-}}
    g(x) = \infty $ ，$ f,g $ 在 $ [a,b) $ 非负且对任意 $ \xi < b $ 都有在 $ [a,\xi] $ 上可积，$ f\leq g $ ，
    则\begin{enumerate}
        \item 若 $ \int_a^{b}g(x)\mathrm{d}x $ 收敛，则$ \int_a^{b}f(x)\mathrm{d}x $收敛；
        \item 若 $ \int_a^{b}f(x)\mathrm{d}x $ 发散，则$ \int_a^{b}g(x)\mathrm{d}x $发散。
    \end{enumerate}
\end{Theo}

\begin{Theo}[比较判别法]

    设 $ {\displaystyle\lim_{x\rightarrow b^-}}f(x) = \infty,{\displaystyle\lim_{x\rightarrow b^-}}
    g(x) = \infty $ ，$ f,g $ 在 $ [a,b) $ 非负
    \begin{enumerate}
        \item 若 $ g(x) > 0, x\in[a,b) $ 且 $ \dis {\displaystyle\lim_{x\rightarrow b^-}}
        \dfrac{f(x)}{g(x)} = k \neq 0 $ ，则 $ \int_a^{b}f(x)\mathrm{d}x $ 与
        $ \int_a^{b}g(x)\mathrm{d}x $ 的敛散性相同；
        \item 若当 $ x\rightarrow b^- $时，$ f\sim g $，则 $ \int_a^{b}f(x)\mathrm{d}x $ 与
        $ \int_a^{b}g(x)\mathrm{d}x $ 的敛散性相同。
    \end{enumerate}
\end{Theo}

\begin{Theo}[绝对收敛必收敛]

    若 $ {\displaystyle\lim_{x\rightarrow b^-}}f(x)=\infty $ 且 $ \dis \int_a^{b}|f(x)|\mathrm{d}x $ 收敛，则
    $ \dis \int_a^{b}f(x)\mathrm{d}x $ 收敛。
\end{Theo}

\section{定积分的几何应用}

主要应用于求面积、体积、弧长、侧面积。

\subsection{平面图形求面积}

\sssubsection{直角坐标系下平面图形求面积}

对于由 $ x = x_1, x = x_2,y = f_1(x), y = f_2(x) (x_1 < x_2,\forall x \in [x_1,x_2]f_1(x) > f_2(x)) $ 
围成的图形，有 $$
    S = \int_a^b[f_1(x)-f_2(x)]\mathrm{d}x
$$ 

对于由 $ y = y_1, y = y_2,x = f_1(y), x = f_2(y) (y_1 < y_2,\forall y \in [y_1,y_2]f_1(y) > f_2(y)) $ 
围成的图形，有 $$
    S = \int_a^b[f_1(y)-f_2(y)]\mathrm{d}y
$$ 

\sssubsection{极坐标下的平面图形求面积}

$$
    S = \int_{\theta_1}^{\theta_2}\textcolor{red}{\dfrac{1}{2}r^2}\mathrm{d}\theta
$$ 

\sssubsection{由参数方程表示的平面图形求面积}

对由 $ \begin{cases}
    x = x(t) \\ y = y(t)
\end{cases} (\alpha\leq t\leq \beta) $ 表示的平面图形，其面积为$\dis\int_\alpha^\beta |y(t)x'(t)|\mathrm{d}t. $

\sssubsection{重要的平面曲线}

重要的平面曲线如下。
\begin{itemize}
    \item 心型线 $ r = a(1+\cos\theta) ; $
    \item 摆线（旋轮线）；
    \item 星型线；
    \item 双扭线；
    \item 阿基米德螺线 $ r = a\theta ; $
    \item 对数螺线；
    \item 三叶玫瑰线 $ r = \sin(3\theta),(0\leq\theta\leq \dfrac{\pi}{3}) ; $
\end{itemize}

\subsection{旋转体求体积}

\sssubsection{薄片法}

绕 $ x $ 轴旋转时，有 $ \mathrm{d}V_x = \pi y^2 \mathrm{d}x $ ，则体积为$$
    V = \int_a^b \pi f(x)^2 \mathrm{d}x
$$ 

\sssubsection{柱壳法}

绕 $ y $ 轴旋转时，有 $ \mathrm{d}V_y = 2\pi xy \mathrm{d} x $ ，则体积为$$
    V = \int_a^b 2\pi \left|xf(x)\right|\mathrm{d}x
$$ 

\subsection{平均值}

设 $ f(x) $ 在区间 $ [a,b] $ 连续，则 $ f(x) $ 在 $ [a,b] $ 的平均值为
$$
    \dfrac{\dis \int_a^b f(x)\mathrm{d}x}{b-a}
$$ 

\subsection{平行截面面积已知的立体体积}

$ \dis V = \int_a^bA(x)\mathrm{d}x,a<b $ ，其中 $ A(x) $ 是待求体积立体截面面积。

\subsection{平面曲线段的弧长}

\sssubsection{直角坐标}

对曲线段 $ y = f(x),x\in[a,b] $ ，设 $ f(x) $ 有连续导数，则给定平面曲线段弧长元素
和弧长分别为
$$ \dis \mathrm{d}s = \dsqrt{1+f'^2(x)}\mathrm{d}x;
s = \int_a^b \dsqrt{1+f'^2(x)}\mathrm{d}x. $$

\sssubsection{参数方程}

若曲线能表示为 $ x = x(t),y=y(t),t\in[\alpha,\beta] $，且其在$ (\alpha,\beta) $ 内
有连续导数，则给定平面曲线段弧长元素和弧长分别为

$$ \dis \mathrm{d}s = \dsqrt[]{y'^2(t)+x'^2(t)}\mathrm{d}t;
s = \int_a^b \dsqrt[]{[y^2(t)]^2+[x'(t)]^2}\mathrm{d}t. $$

\sssubsection{极坐标}


若曲线能表示为 $ \rho = \rho(\theta),\theta\in[\theta_1,\theta_2], $
则给定平面曲线段弧长元素和弧长分别为

$$ \dis \mathrm{d}s = \dsqrt[]{\rho^2(\theta)+\rho'^2(\theta)}\mathrm{d}\theta ;
s = \int_a^b \dsqrt[]{\rho^2(\theta)+[\rho'(\theta)]^2}\mathrm{d}\theta. $$

\subsection{旋转曲面求侧面积}

曲线$ y = f(x),x\in(a,b) $ 弧绕 $ x $ 轴旋转所得曲面表面积为
$ \dis S = 2\pi\int_a^b|f(x)|\overbrace{\dsqrt[]{1+[f'(x)^2]}\mathrm{d}x}^{\mathrm{d}s} $ 

曲线$ y = f(x),x\in(a > 0,b) $ 弧绕 $ y $ 轴旋转所得曲面表面积为
$ \dis S = 2\pi\int_a^bx\cdot\overbrace{\dsqrt[]{1+[f'(x)^2]}\mathrm{d}x}^{\mathrm{d}s} $ 

当旋转曲面使用极坐标或参数方程表示时，将 $ \mathrm{d}s $ 改为相对应的即可。

\section{定积分的物理应用}
定积分的物理应用主要有以下几种。

\begin{itemize}
    \item 做功 $ W = FS; $ 
    \item 受力 \begin{itemize}
        \item 万有引力 $ F = \dfrac{GMm}{r^2}; $ 
        \item 液体压力 $ F = PS = \rho ghS; $ 
    \end{itemize}
\end{itemize}


\chapter{线性方程组}

\begin{equation*}
    \begin{aligned}
        \begin{cases}
            \textrm{齐次方程组} AX = 0\begin{cases}
                \textrm{表示形式}\begin{cases}
                    \textrm{矩阵}\\ \textrm{向量}
                \end{cases}\\ 
                \textrm{有解的判定}\\ 
                \textrm{解的性质}\begin{cases}
                    \textrm{通解}\\ \textrm{基解}\\ \textrm{组合}
                \end{cases}\\ 
                \textrm{求解方法}
            \end{cases}\\ 
            \textrm{非齐次方程组} AX = b\begin{cases}
                \textrm{有解判定} \\ \textrm{解的性质} \\ \textrm{解的结构 - 齐通+非齐特}
            \end{cases}\\ 
            \textrm{同解与公共解}\begin{cases}
                \textrm{通解 - 解集相同}\\ \textrm{公共解 - 解集有交集}
            \end{cases}
        \end{cases}
    \end{aligned}
\end{equation*}

\section{解的判定}

对于齐次方程 $ AX = 0, \begin{cases}
    \textrm{只有零解}\Rightarrow r(A) = n\\ 
    \textrm{有非零解}\Rightarrow r(A) < n
\end{cases}. $ 

其基础解系
\begin{itemize}
    \item 为 $ AX = 0 $ 的解；
    \item 线性无关；
    \item 数量为 $ n - r(A) $ 个。
\end{itemize}

对于非齐次方程 $ AX = b,\begin{cases}
    \textrm{无解}\Rightarrow r(A) \neq r(\bar A)\\ 
    \textrm{有解}\Rightarrow r(A) = r(\bar A)
    \begin{cases}
         = n \Rightarrow \textrm{解唯一}\\ 
         < n \Rightarrow \textrm{解不唯一（无穷）}
    \end{cases}
\end{cases} $ 

\section{求通解}

对于具体的数字矩阵，通过行变换或高斯消元做。

对抽象矩阵，
\begin{itemize}
    \item 通解为齐通 + 非齐特；
    \item 齐解任意组合仍然为齐解；
    \item 若 $ \xi_1,\cdots,\xi_t $ 均为 $ AX = b $ 的解，则
    $ \dis \sum k_i\xi_i $ 为 $ \begin{cases}
        \textrm{齐解，} \sum k = 0 \\
        \textrm{非齐解，} \sum k = 1 \\
    \end{cases} $ 
\end{itemize}
\chapter{二重积分}

\section{交换积分次序}

$ \begin{cases}
    \textrm{直角 - 穿针法}\\
    \textrm{极坐标}\begin{cases}
        xy\textrm{坐标系 - 作圆弧}\\ 
        \theta r\textrm{坐标系 - 穿针法(当成直角坐标)}
    \end{cases}
\end{cases} $ 

二重变限积分求极限时，
$ \begin{cases}
    \textrm{洛必达结合交换积分次序}\\\textrm{二重积分中值定理(非零因子)}
\end{cases} $ 

\sssubsection{例题T6.2.2}

\begin{equation*}
    \begin{aligned}
        {\displaystyle\lim_{t\rightarrow 0^+}}\dfrac{1}{t^3}\int_0^t\mathrm{d}x
        \int_{x^2}^{t^2}\arctan(\cos(2x+3\dsqrt{y}))\mathrm{d}y &= 
        {\displaystyle\lim_{t\rightarrow 0^+}}\dfrac{1}{t^3}\iint\limits_{D}
        \arctan(\cos(2x+3\dsqrt{y}))\mathrm{d}x\mathrm{d}y
        \\&= {\displaystyle\lim_{t\rightarrow 0^+}}\dfrac{1}{t^3}
        \arctan(\cos(2\xi+3\dsqrt{\eta}))\cdot S_D \,(\xi,\eta \in D)
        \\&= {\displaystyle\lim_{t\rightarrow 0^+}}\dfrac{1}{t^3}
        \arctan(\cos(2\xi+3\dsqrt{\eta}))\cdot\dfrac{2}{3}t^3
        \\&= \dfrac{\pi}{4}\cdot\dfrac{2}{3} = \dfrac{\pi}{6}.
    \end{aligned}
\end{equation*}

\sssubsection{例题T6.3.2}

\begin{equation*}
    \begin{aligned}
        \int_0^\frac{\pi}{2}\dfrac{1}{\sqrt{x}}\mathrm{d}x\int_{\sqrt{x}}^{\sqrt{\frac{\pi}{2}}}
        \dfrac{1}{1+(\tan y^2)^{\sqrt 2}}\mathrm{d}y 
        &= \int_0^{\sqrt{\frac{\pi}{2}}}\dfrac{1}{1+(\tan y^2)^{\sqrt 2}}
        \mathrm{d}y\int_0^{y^2}\dfrac{1}{\sqrt x}\mathrm{d}x
        \\&= \int_0^{\sqrt{\frac{\pi}{2}}}\dfrac{2y}{1+(\tan y^2)^{\sqrt 2}}\mathrm{d}y
        \\&= \int_0^{\frac{\pi}{2}}\dfrac{\mathrm{d}t}{1+(\tan t)^{\sqrt 2}}
        \\&= \int_0^{\frac{\pi}{2}}\dfrac{(\cos t)^{\sqrt 2}\mathrm{d}t}{(\sin t)^{\sqrt 2}+(\cos t)^{\sqrt 2}} 
        \\\textrm{区间再现}&=\dfrac{1}{2}\cdot\dfrac{\pi}{2} = \dfrac{\pi}{4}.
    \end{aligned}
\end{equation*}

\sssubsection{例题T6.3.4}

\begin{equation*}
    \begin{aligned}
        \int_0^1\dfrac{x^b-x^a}{\ln x}\mathrm{d}x &= 
        \int_0^1 \mathrm{d}x\int_a^b x^y\mathrm{d}y \\ 
        &= \int_a^b \mathrm{d}y \int_0^1 x^y\mathrm{d}x
        \\&= \int_a^b \dfrac{\mathrm{d}y}{1+y} = \ln\dfrac{1+b}{1+a}.
    \end{aligned}
\end{equation*}

\sssubsection{例题T6.3.5}

此处以 $ \theta - r $ 坐标系交换积分次序。

\begin{equation*}
    \begin{aligned}
        \int_0^{2\pi}(\theta^2-1)\mathrm{d}\theta\int_{\frac{\theta}{2}}^\pi e^{r^2}\mathrm{d}r &=
        \int_0^\pi e^{r^2}\mathrm{d}r\int_0^{2r}(\theta^2-1)\mathrm{d}\theta \\ 
        &= \dfrac{8}{3}\int_0^\pi r^3e^{r^2}\mathrm{d}r - 2\int_0^\pi re^{r^2}\mathrm{d}r
        \\ &= \dfrac{1}{3}(4\pi^2-7)e^{\pi^2}+\dfrac{7}{3}.
    \end{aligned}
\end{equation*}

\section{二重积分的计算}

考虑
$ \begin{cases}
    \textrm{直角坐标系}\\
    \textrm{极坐标变换}\\
    \textrm{奇偶性}\\
    \textrm{轮换对称性}\\
    \textrm{形心公式}\\
    \textrm{平移变换}\\
    \textrm{雅可比行列式}
\end{cases} $ 

\sssubsection{被积函数为齐次函数 $ f(\dfrac{y}{x}) $ T6.6}

\begin{enumerate}
    \item[\textbf{例题}] 设有区域 $ D = \{(x,y)|x|+|y|\leq 1\}, $ 
    计算 $ \dis \iint\limits_D e^{\dfrac{|y|}{|x|+|y|}}\mathrm{d}x\mathrm{d}y. $ 
    \item[\textrm{\textbf{方法}}] 可以知道，
    \begin{equation*}
        \begin{aligned}
            \textrm{原积分} &= 4\iint\limits_{D_1}e^{\dfrac{y}{x+y}}\mathrm{d}x\mathrm{d}y \\ 
            &= 4\int_0^{\frac{\pi}{2}}\mathrm{d}\theta\int_0^{\frac{1}{\sin \theta + \cos \theta}}
            e^{\frac{\sin \theta}{\sin \theta + \cos \theta}}r\mathrm{d}r
            \\&= 2\int_0^{\frac{\pi}{2}} e^{\frac{\sin \theta}{\sin \theta + \cos \theta}} 
            \dfrac{1}{(\cos \theta + \sin \theta)^2}\mathrm{d}\theta
            \\&= 2\int_0^{\frac{\pi}{2}} e^{\frac{\sin \theta}{\sin \theta + \cos \theta}} 
            \mathrm{d}\dfrac{\sin \theta}{\sin \theta + \cos \theta}
            \\&= 2e^{\frac{\sin \theta}{\sin \theta + \cos \theta}}
            \Big|_0^\frac{\pi}{2} = 2(e-1).
        \end{aligned}
    \end{equation*}
\end{enumerate}

事实上，对于包含齐次函数 $ f(\dfrac{y}{x}) $ 的二重积分，都应该考虑利用极坐标变出 
$ \mathrm{d}\dfrac{A\sin\theta+B\cos\theta}{\sin\theta+\cos\theta}. $

对于 $ f(\dfrac{y^2}{x^2}),f(\dfrac{ax+by}{x+y}) $ 也是类似。

\sssubsection{极坐标压轴题 - T6.7}

\begin{enumerate}
    \item[\textbf{例题}] 设 $ f(x,y) $ 在区域 $ D = \{(x,y)|x^2+y^2\leq 1\} $ 
    上有一阶连续偏导数，在 $ D $ 的边界上 $ f(x,y) $ 恒为零，求
    $$
        {\displaystyle\lim_{\varepsilon\rightarrow 0^+}}\iint\limits_{\varepsilon^2\leq x^2+y^2\leq 1}
        \dfrac{xf'_x+yf'_y}{x^2+y^2}\mathrm{d}x\mathrm{d}y.
    $$
    \item[\textbf{方法}] 引入极坐标，则有
    \begin{equation*}
        \begin{aligned}
            &\dfrac{\partial f}{\partial \rho} = 
            \dfrac{\partial f}{\partial x}\dfrac{\partial x}{\partial \rho}
            + \dfrac{\partial f}{\partial y}\dfrac{\partial y}{\partial \rho}
            = f'_x\cos\theta + f'_y\sin\theta \\ 
            &\dfrac{\partial f}{\partial \theta} = 
            \dfrac{\partial f}{\partial x}\dfrac{\partial x}{\partial \theta}
            + \dfrac{\partial f}{\partial y}\dfrac{\partial y}{\partial \theta}
            = -f'_x(\rho \sin\theta)+f'_y(\rho \cos\theta)\\ 
            \Rightarrow & \rho \dfrac{\partial f}{\partial \rho} = 
            \rho\left[f'_x\cos\theta + f'_y\sin\theta\right]
            =xf'_x+yf'_y.
        \end{aligned}
    \end{equation*}
    那么，对于题设，有
    \begin{equation*}
        \begin{aligned}
            \textrm{原极限}&= {\displaystyle\lim_{\varepsilon\rightarrow 0^+}}
            \int_0^{2\pi}\mathrm{d}\theta\int_{\varepsilon}^1
            \dfrac{\partial f}{\partial \rho}\mathrm{d}\rho \\ 
            &= {\displaystyle\lim_{\varepsilon\rightarrow 0^+}}
            \int_0^{2\pi}f(\rho\cos\theta,\rho\sin\theta)\Big|_\varepsilon^1\mathrm{d}\theta\\
            &= -{\displaystyle\lim_{\varepsilon\rightarrow 0^+}}
            \int_0^{2\pi}f(\varepsilon\cos\theta,\varepsilon\sin\theta)\mathrm{d}\theta \\ 
            &= -\int_0^{2\pi}f(0,0)\mathrm{d}\theta = -2\pi f(0,0).
        \end{aligned}
    \end{equation*}
\end{enumerate}

\sssubsection{极坐标逆问题 - T6.8}

在极坐标不好做的情况下，将 $ (\rho,\theta) $ 转换回 $ (x,y) $ 即可。

\newpage

\sssubsection{奇偶性 - T6.9-10}

\begin{itemize}
    \item \textbf{奇偶性的推广}
    
    设 $ f(x,y) $ 在区域 $ D = D_1\bigcup D_2 $ 上连续，其中 $ D_1,D_2 $ 关于原点对称，则
    $$
        \iint\limits_Df(x,y)\mathrm{d}x\mathrm{d}y
        =\begin{cases}
            2\dis\iint\limits_{D_1}f(x,y)\mathrm{d}x\mathrm{d}y,\quad{} f(-x,-y) = f(x,y)\\
            0,\quad{} f(-x,-y) = -f(x,y)\\
        \end{cases}
    $$
\end{itemize}

\begin{enumerate}
    \item[\textbf{例题}]
    设区域 $ D $ 由曲线 $ y = x^3,y = x $ 围成。求
    $$
        \iint\limits_D\left[e^{x^2}+\sin(x+y)\right]\mathrm{d}x\mathrm{d}y.
    $$
    \item[\textbf{方法}] 
    由奇偶性，可以知道
    \begin{equation*}
        \begin{aligned}
            \textrm{原积分}&=2\int_0^1\mathrm{d}x\int_{x^3}^x e^{x^2}\mathrm{d}y\\ 
            &= 2\int_0^1e^{x^2}(x-x^3)\mathrm{d}x 
            \\&= \int_0^1(1-x^2)e^{x^2}\mathrm{d}x^2 = e-2.
        \end{aligned}
    \end{equation*}
\end{enumerate}

\begin{enumerate}
    \item[\textbf{例题}]
    对 $ D = \{(x,y)|x^2+y^2\leq 1, x+y\geq 0\}, $ 求
    $$
        \iint\limits_D\dfrac{1+2x^2+xy^2}{1+x^2+y^2}\mathrm{d}x\mathrm{d}y.
    $$
    \item[\textbf{方法}] 
    原积分等于 $ \dis \iint\limits_D\dfrac{1+2x^2}{1+x^2+y^2}\mathrm{d}x\mathrm{d}y+
    \iint\limits_D\dfrac{xy^2}{1+x^2+y^2}\mathrm{d}x\mathrm{d}y. $ 
    前一部分利用轮换对称性。对后一部分，以 $ y = x $ 将 $ D $ 
    划分为 $ D_1,D_2, $ 以 $ y = 0 $ 将 $ D_1 $ 划分为 $ D_{11},D_{12}. $ 
    那么，
    \begin{equation*}
        \begin{aligned}
            \textrm{后半部分积分} &= \iint\limits_{D_1}+\iint\limits_{D_2}
            \dfrac{xy^2}{1+x^2+y^2}\mathrm{d}x\mathrm{d}y\\ 
            &= 2\iint\limits_{D_12}\dfrac{xy^2}{1+x^2+y^2}\mathrm{d}x\mathrm{d}y\\ 
            &= 2\int_{0}^\frac{\pi}{4}\mathrm{d}\theta\int_0^1
            \dfrac{r^4\sin^2\theta\cos\theta}{1+r^2}\mathrm{d}r = \dfrac{1}{3\sqrt 2}
            \left(\dfrac{\pi}{4}-\dfrac{2}{3}\right).
        \end{aligned}
    \end{equation*}
\end{enumerate}

\sssubsection{轮换对称性的推广 - T6.11}

\begin{itemize}
    \item \textbf{轮换对称性的推广}
    
    $$
        \iint\limits_Df(x,y)\mathrm{d}\sigma = 
        \iint\limits_{D'}f(y,x)\mathrm{d}\sigma
    $$
    其中 $ D' = \{(y,x)|(x,y)\in D\}. $ 
\end{itemize}
\chapter{无穷级数}

\section{数项级数}

\subsection{数项级数概念}

\begin{Def}[数项级数]

    无穷多个数 $ u_1,u_2,\dots,u_n,\dots $ 按照一定的顺序相加得到的表达式
     $ \sum_{n=1}^\infty u_n $ 称为数项级数，简称级数，其中 $ u_n $ 称为通项。
\end{Def}

级数是无数个的，也是有序的。

\begin{Def}[级数收敛]

    称 $ S_n = \sum_{k=1}^n u_k $ 为级数的部分和，数列 $ \{S_n\} $ 即为部分和数列。

    若 $ {\displaystyle\lim_{n\rightarrow \infty}}S_n\xlongequal{\exists}S $ ，
    则称级数 $\dis \sum_{n=1}^\infty u_n $ 收敛，且和为 $ S $ ;反之，若极限不存在，则称其
    为发散的，发散级数没有和的概念。
\end{Def}

以下是几个重要级数及其敛散性。

\begin{itemize}
    \item 几何级数 - $\dis \sum_{n=0}^\infty ar^n $ ，当 $ |r|<1 $ 时，其收敛于 $ \dfrac{a}{1-r} $ ;
    若 $ |r|\geq1 $ ，则其发散。
    \item p级数 - $\dis \sum_{n=1}^\infty \dfrac{1}{n^p} $ ，若 $ p > 1 $ ，其收敛；若 $ p\leq 1 $ ，其发散。
\end{itemize}

事实上，有 $\dis \sum_{n=1}^\infty \dfrac{1}{n^2} = \dfrac{\pi^2}{6}. $ 

\section{级数基本性质及其收敛必要条件}

级数有以下性质。

\begin{itemize}
    \item 若 $ \dis \sum_{n=1}^{\infty}u_n,\sum_{n=1}^{\infty}v_n $ 都收敛，则
    $ \dis \sum_{n = 1}^{\infty}(u_n\pm v_n) $ 也收敛，且有
    $ \dis \sum_{n = 1}^{\infty}(u_n\pm v_n) = \sum_{n=1}^{\infty}u_n\pm\sum_{n=1}^{\infty}v_n $ .
    \item 在级数中增加、减少、改变有限项不会影响级数敛散性。
    \item 收敛级数具有结合律，即可以对级数的项任意地添加括号，新级数仍然收敛于同一个和。
    注意，去括号不成立。
    \item 对 $ k \neq 0 $ 有 $ \dis \sum_{n=1}^{\infty}ku_n $ 与  $ \dis \sum_{n=1}^{\infty}u_n $
    具有相同的敛散性。
    \item 级数  $ \dis \sum_{n=1}^{\infty}u_n $ 收敛的必要条件是 $ {\displaystyle\lim_{n\rightarrow \infty}}u_n = 0 $ .
\end{itemize}

\section{正项级数及其敛散性判别法}

\begin{Def}[正项级数]

    若 $ u_n\geq 0 $ ，则有 $\dis \sum_{n=1}^{\infty}u_n $ 称为正项级数。
\end{Def}

\begin{Theo}[正项级数收敛基本定理]

    正项级数 $\dis \sum_{n=1}^{\infty}u_n $ 收敛 $ \Leftrightarrow  S_n $ 有上界。
\end{Theo}

正项级数有以下几种判别法。

\begin{Theo}[直接比较法]

    设当 $ n\geq N $ 时，有 $ 0\leq u_n\leq v_n $ 总是成立，则\begin{itemize}
        \item 若 $\dis \sum_{n=1}^{\infty}v_n $ 收敛，则 $\dis \sum_{n=1}^{\infty}u_n $ 收敛；
        \item 若 $\dis \sum_{n=1}^{\infty}u_n $ 发散，则 $\dis \sum_{n=1}^{\infty}v_n $ 发散。
    \end{itemize}
\end{Theo}

\begin{Theo}[间接比较法]

    设 $ u_n\geq 0, v_n>0,n = 1,2,3,\dots $，且有 $ {\displaystyle\lim_{n\rightarrow \infty}}
    \dfrac{u_n}{v_n} = A  $ ，则\begin{itemize}
        \item 若 $ 0<A<+\infty $ ，两级数具有相同的敛散性；
        \item 若 $ A = 0 $ ，则若 $\dis \sum_{n=1}^{\infty}v_n $ 收敛，有 $\dis \sum_{n=1}^{\infty}u_n $ 收敛；
        \item 若 $ A = +\infty $ ，若 $\dis \sum_{n=1}^{\infty}u_n $ 收敛，则 $\dis \sum_{n=1}^{\infty}v_n $ 收敛。
    \end{itemize}
\end{Theo}

\begin{Theo}[比值判别法]

    设 $ u_n > 0 $ ，有 $ {\displaystyle\lim_{n\rightarrow \infty}}\dfrac{u_{n+1}}{u_{n}} = \rho $ ，
    则\begin{itemize}
        \item 若 $ \rho > 1 $ ，有 $\dis \sum_{n=1}^{\infty}u_n $ 收敛；
        \item 若 $ \rho < 1 $ ，则其发散；
        \item 若 $ \rho = 1 $ ，该判别法失效。
    \end{itemize}
\end{Theo}

\begin{Theo}[根值判别法]

    设 $ u_n \leq 0 $ ，有 $ {\displaystyle\lim_{n\rightarrow \infty}}\sqrt[n]{u_n} = \rho $ ，
    则\begin{itemize}
        \item 若 $ \rho > 1 $ 甚至为$ +\infty $ 时，有 $\dis \sum_{n=1}^{\infty}u_n $ 收敛；
        \item 若 $ \rho < 1 $ ，则其发散；
        \item 若 $ \rho = 1 $ ，该判别法失效。
    \end{itemize}
\end{Theo}

\begin{Theo}[积分准则]

    若函数在 $ [1,+\infty) $ 上非负、连续、递减，则级数 $\dis \sum_{n=1}^{\infty}f(n) $ 与
    积分 $ \dis \int_1^{+\infty}f(x)\mathrm{d}x $ 具有相同的敛散性。
\end{Theo}

注意，判别正项级数敛散性时，\begin{equation*}
    \begin{matrix}
        \textrm{间接比较法} \\\Downarrow\\ 
        \textrm{比值法}\\\Downarrow\\
         \textrm{根值法}\\\Downarrow\\
          \textrm{积分准则}\\\Downarrow\\
        \textrm{直接比较法}\\\Downarrow\\
         \textrm{定义·部分和数列}\\\Downarrow\\
          \textrm{收敛级数定义}
    \end{matrix}
\end{equation*}

此外，若通项为积分或者是抽象的，从直接比较法开始应用，试图放缩被积函数以去除积分号。

\section{任意项级数}

\subsection{交错级数及其敛散性判别法}

\begin{Def}[交错级数]

    若 $\dis u_n>0,\sum_{n=1}^{\infty}(-1)^{n+1}u_n $ 称为交错级数。
\end{Def}

\begin{Theo}[莱布尼茨判别法]

    交错级数 $\dis \sum_{n=1}^{\infty}(-1)^{n+1}u_n $ 满足
    \begin{itemize}
        \item $ u_{n+1}\leq u_n $ 
        \item $ {\displaystyle\lim_{n\rightarrow \infty}}u_n = 0 $ 
    \end{itemize}
    则该交错级数收敛，且 $\dis 0\leq \sum_{n=1}^{\infty}(-1)^{n+1}u_n\leq u_1 $ .
\end{Theo}

\subsection{条件收敛与绝对收敛}

\begin{Def}[条件收敛与绝对收敛]

    若 $ \dis \sum_{n=1}^{\infty}|u_n| $ 收敛，则称 $ \dis \sum_{n=1}^{\infty}u_n $ 绝对收敛。

    若 $ \dis \sum_{n=1}^{\infty}u_n $ 收敛而 $ \dis \sum_{n=1}^{\infty}|u_n| $ 发散，则
    $ \dis \sum_{n=1}^{\infty}u_n $ 称为条件收敛。
\end{Def}

绝对收敛的级数有交换律，而条件收敛的级数没有。

\begin{Theo}[条件收敛必收敛]

    $ \dis \sum_{n=1}^{\infty}|u_n|$ 收敛 $\dis \Rightarrow \sum_{n=1}^{\infty}u_n $ 收敛。
\end{Theo}

\section{幂级数}

\subsection{函数项级数}

\begin{Def}[函数项级数]

    设 $ u_n(x),n=1,2,\dots $ 定义在 $ I $ 上，则 $ \dis \sum_{n=1}^{\infty}u_n(x) $ 
    称为区间 $ I $ 上的函数项级数。
\end{Def}

\begin{Def}[收敛点与收敛域]

    对 $ x_0\in I $ ，若常数项级数 $\dis \sum_{n=1}^{\infty}u_n(x_0) $ 收敛，则称 $ x_0 $ 为函数项级数
    $ \dis \sum_{n=1}^{\infty}u_n(x) $ 的收敛点；反之，若其发散，则为发散点。

    函数项级数 $ \dis \sum_{n=1}^{\infty}u_n(x) $ 的全部收敛点构成的集合是收敛域，全部发散点构成的集合
    是发散域。
\end{Def}

\begin{Def}[和函数]

    和函数 $\dis S(x) =: \sum_{n=1}^{\infty}u_n(x) $ ，其定义域为收敛域。
\end{Def}

\subsection{幂级数}

\begin{Def}[幂级数]

    形如$ \dis \sum_{n=1}^\infty a_n(x-x_0)^n $ 的级数称为 $ (x-x_0) $ 的幂级数，常数 $ a_n $ 称为
    幂级数的系数。
\end{Def}

显然，当 $ x_0 = 0 $ 时，有 $ \dis \sum_{n=0}^{\infty}a_nx^n $ 称为 $ x $ 的幂级数。

\begin{Theo}[阿贝尔定理]

    若幂级数 $\dis \sum u_n(x) $ 在 $ x = x_0 $ 处收敛，则对任意 $ x $ 满足 $ |x| < |x_0| $ ,
    都有$ \dis \sum u_n(x) $ 绝对收敛；
    若在 $ x = x_0 $ 发散，则对任意 $ x $ 满足 $ |x|>|x_0| $ 都有$\dis \sum u_n(x) $ 发散。
\end{Theo}

\subsection{幂级数收敛半径及其求法}

幂级数的收敛半径 $ R $ 可以由以下公式求出。

\begin{itemize}
    \item $ \dis R = {\displaystyle\lim_{n\rightarrow \infty}}\left|\dfrac{a_n}{a_{n+1}}\right| $ 
    \item $ \dis R = {\displaystyle\lim_{n\rightarrow \infty}}\dfrac{1}{\dsqrt[n]{|a_n|}} $ 
\end{itemize}

注意，公式仅对标准幂级数成立，即要求幂函数的次数逐次增长；

收敛区间的中心为 $ x_0 $. 以 $ x = 0 $ 为例子，若存在收敛半径 $ R\geq 0 $ ，则
$ x\in (-R,R) $ 时级数绝对收敛，$ x\in R\big\backslash(-R,R) $ 时级数发散。
若级数有条件收敛点，则其只能在 $ \pm R $ 处，因而最多有两个。

在求收敛半径后，必须写出收敛区间，然后判断端点处敛散性。

\subsection{幂级数的性质}

\sssubsection{幂级数的加减运算}

设 $ \dis \sum_{n=1}^\infty a_nx^n = f(x),|x|< R_1;\sum_{n=1}^\infty b_nx^n = g(x), |x|< R_2 $，则有
$ \dis \sum_{n=1}^\infty (a_n\pm b_n)x^n = f(x)\pm g(x), |x|<\min(R_1,R_2)  $ .

\sssubsection{幂级数的逐项求导与逐项积分}

设幂级数 $ \dis \sum_{n=0}^\infty a_nx^n $ 的收敛半径 $ R>0 $ 且有和函数 $ \dis S(x) = \sum_{n=1}^\infty a_nx^n $ ，则
成立以下性质。
\begin{itemize}
    \item $ S(x) $ 在 $ \overset{\textrm{收敛区间}}{(-R,R)} $ 内可导，且有逐项求导公式$$
        S'(x) = \left(\sum_{n=0}^\infty a_nx^n\right)' = \sum_{n=0}^\infty (a_nx^n)' = \sum_{n=0}^\infty na_nx^{n-1}
    $$ 
    求导后幂级数收敛半径不变，因而 $ S(x) $ 在 $ (-R,R) $ 内有任意阶导数；
    \item $ S(x) $ 在$ \overset{\textrm{收敛区间}}{(-R,R)} $内有逐项积分公式$$
        \int_0^x S(t)\mathrm{d}t = \sum_{n=0}^\infty\int_0^x a_nt^n\mathrm{d}t = \sum_{n=0}^\infty \dfrac{a_n}{n+1}x^{n+1}
    $$ 
    积分后收敛半径不变；
    \item 幂级数的和函数 $ S(x) $ 在收敛域上连续。
\end{itemize}

\sssubsection{单位幂级数}

$ \dis \sum_{n=0}^\infty x^n $ 称为单位幂级数,显然其收敛半径是1,且收敛区间内，有 $ \dis \sum x^n = \dfrac{1}{1-x} $ .

$ \dis \sum_{n=0}^\infty (-1)^{n+1}x^n $ 收敛半径也是1, 且收敛区间内，有 $ \dis \sum x^n = \dfrac{1}{1+x} $ .


\section{函数展开成幂级数}

\subsection{概念}

设 $ f(x) $ 是给定的函数，若存在幂级数 $ \dis \sum_{n=0}^\infty a_n(x-x_0)^n $ 使得任意 $ x\in I $ 
级数 $ \dis \sum_{n=0}^\infty a_n(x-x_0)^n $ 都收敛到 $ f(x) $ ，即 $\dis \forall x\in I, f(x) = \sum_{n=0}^\infty a_n(x-x_0)^n $ ,
则称 $ f(x) $ 在区间 $ I $ 上可展开为 $ x - x_0 $ 的幂级数，称级数 $ \dis \sum_{n=0}^\infty a_n(x-x_0)^n $ 为$ f(x) $ 的幂级数。

\begin{Def}[幂级数展开唯一性]

    若 $ f(x) $ 在 $ x = x_0 $ 可以展开为幂级数 $ \dis \sum_{n=0}^\infty a_n(x-x_0)^n $ ，则该展开式是唯一的，且为
    泰勒级数，即有 $ \dis f(x) = \sum_{n=0}^\infty a_n(x-x_0)^n = \sum_{n=0}^\infty \dfrac{f^{(n)(x_0)}}{n!}(x-x_0)^n$
\end{Def}

\subsection{函数展开成幂级数的条件及其形式}

\sssubsection{泰勒级数与马克劳林级数的定义}

设函数 $ f(x) $ 在 $ x_0 $ 的某一领域 $ |x-x_0|<\sigma $ 内具有任意阶导数，则级数
$ \dis \sum_{n=0}^\infty \dfrac{f^{(n)(x_0)}}{n!}(x-x_0)^n $ 称为函数 $ f(x) $ 在 $ x = x_0 $ 处的泰勒级数。

特别地，若 $ x_0 = 0 $ ，称级数为$ f(x) $ 的马克劳林级数。

\sssubsection{函数展开成为幂级数的充要条件}

设 $ f(x) $ 在 $ |x-x_0|<R $ 内有任意阶导数，其泰勒公式为$$
    f(x) = \sum_{i=0}^n \dfrac{f^{(i)}(x_0)}{i!}(x-x_0)^i + R_n(x)
$$ 其中$ R_n(x) $为 $ n $ 阶余项，
其拉格朗日型为 $ \dis R_n(x) = \dfrac{f^{(n+1)}[x_0+\theta(x-x_0)]}{(n+1)!}(x-x_0)^{n+1},\theta\in(0,1) $ ,
则 $ \dis f(x) = \sum_{n=0}^\infty \dfrac{f^{(n)}(x_0)}{ni!}(x-x_0)^n,|x-x_0|<R $ 的充要条件为
$ {\displaystyle\lim_{n\rightarrow \infty}}R_n(x) = 0,|x-x_0|<R $ .

特别地，当 $ x_0 = 0 $ 时得到函数展开为马克劳林级数的充要条件。

\subsection{函数展开成为幂级数的方法}

\begin{itemize}
    \item 公式法 $$
     \dis f(x) = \sum_{n=0}^\infty \dfrac{f^{(n)}(x_0)}{ni!}(x-x_0)^n,|x-x_0|<R 
    $$
    \item 变量替换；
    \item 逐项求导；
    \item 逐项积分；
    \item 恒等变形、变量代换等。
\end{itemize}

幂级数有8个常用的展开式。

\begin{itemize}
    \item $ \dis e^x = \sum_{n=0}^{\infty}\dfrac{1}{n!}x^n, |x|<\infty $ 
    \item $ \dis \sin x = \sum_{n=0}^{\infty}(-1)^n\dfrac{x^{2n+1}}{(2n+1)!},|x|<\infty $ 
    \item $ \dis \cos x = \sum_{n=0}^{\infty}(-1)^n\dfrac{x^{2n}}{(2n)!},|x|<\infty $ 
    \item $ \dis (1+x)^\alpha = 1 + \sum_{n=1}^{\infty}\dfrac{\mathrm{C}_\alpha^n}{n!}x^n, |x|<1 $ 
    \item $ \dis \dfrac{1}{1-x} = \sum_{n=0}^{\infty}x^n, |x|<1 $ 
    \item $ \dis \dfrac{1}{1+x} = \sum_{n=0}^{\infty}(-1)^nx^n, |x|<1 $ 
    \item $ \dis \ln(1+x) = \sum_{n=1}^{\infty} \dfrac{(-1)^{n+1}x^n}{n!}, -1<x\leq 1 $ 
    \item $ \dis \ln(1-x) = \sum_{n=1}^{\infty} \dfrac{x^n}{n!}, -1\leq x< 1 $ 
\end{itemize}

以上都展成马克劳林级数。若要展成泰勒级数，经过适当处理后可以利用马克劳林级数的结果。

展开的结果尽量合成一个级数，级数后要注收敛范围。

可以利用幂级数求数项级数，即将数项级数视为幂级数的 $ x $ 取了定值。

\section{傅里叶级数}

\subsection{傅里叶级数及傅里叶系数}

设函数 $ f(x) $ 在区间 $ [-\pi,\pi] $ 上黎曼可积，则称公式
\begin{itemize}
    \item $ \dis a_k = \dfrac{1}{\pi}\int_{-\pi}^\pi f(x)\cos kx\mathrm{d}x $ 
    \item $ \dis b_k = \dfrac{1}{\pi}\int_{-\pi}^\pi f(x)\sin kx\mathrm{d}x $ 
\end{itemize}

为函数 $ f(x) $ 的傅里叶系数。称以 $ a_k,b_k $ 为系数的三角级数$$
    \dfrac{a_0}{2} + \sum_{n=1}^{\infty}\left(a_n\cos nx + b_n\sin nx\right)
$$ 
为 $ f(x) $ 的傅里叶级数，记为
$$
    f(x)\sim \dfrac{a_0}{2} + \sum_{n=1}^{\infty}\left(a_n\cos nx + b_n\sin nx\right)
$$ 

类似地，可以定义周期为 $ 2l $ 的函数的傅里叶级数，即

\begin{itemize}
    \item $ \dis a_k = \dfrac{1}{l}\int_{-l}^l f(x)\cos \dfrac{k\pi x}{l}\mathrm{d}x $ 
    \item $ \dis b_k = \dfrac{1}{l}\int_{-l}^l f(x)\sin \dfrac{k\pi x}{l}\mathrm{d}x $
    \item $ \dis f(x) \sim \dfrac{a_0}{2} + 
    \sum_{n=1}^{\infty}\left(a_n\cos \dfrac{n\pi x}{l} + b_n\sin \dfrac{n\pi x}{l}\right) $ 
\end{itemize}

\subsection{傅里叶级数的收敛定理}

设 $ f(x) $ 是周期为 $ 2\pi $ 的周期函数，并满足迪利克雷条件，即
\begin{itemize}
    \item 在一个周期内连续或有有限个第一类间断点；
    \item 在一个周期内有有限个极值点，即单调区间个数有限，
\end{itemize}
则$ f(x) $ 的傅里叶级数收敛，且有

\begin{equation*}
    \begin{aligned}
        \dfrac{a_0}{2} + \sum_{n=1}^{\infty}\left(a_n\cos nx + b_n\sin nx\right)
        = \begin{cases}
            f(x),&x\textrm{为$ f(x) $ 的连续点}\\\dfrac{f(x-0)+f(x+0)}{2},& x\textrm{为$ f(x) $ 的间断点}
        \end{cases}
    \end{aligned}
\end{equation*}

\subsection{奇偶函数的傅里叶级数}

若 $ f(x) $ 是奇函数，则有 $ \dis a_n = \dfrac{1}{l}\int_{-l}^l f(x)\cos \dfrac{n\pi x}{l}\mathrm{d}x = 0,
f(x) = \sum_{n=1}^{\infty}b_n\sin \dfrac{n\pi x}{l} $ ，是为正弦级数。

若 $ f(x) $ 是偶函数，则有 $ \dis b_n = \dfrac{1}{l}\int_{-l}^l f(x)\sin \dfrac{n\pi x}{l}\mathrm{d}x = 0,
f(x) =\dfrac{a_0}{2}+  \sum_{n=1}^{\infty}a_n\cos \dfrac{n\pi x}{l} $ ，是为余弦级数。

以上二级数，都有 $ x\in\{f\textrm{的连续点}\} $.

\subsection{有限区间上的函数的傅里叶级数}

将有限区间上的函数展为正弦级数时，将$ f(x) $ 延拓为周期为 $ 2\pi $ 的奇函数 $ F(x) $ .
此时 $ a_n = 0,\dis b_n = \int_{-\pi}^{\pi} f(x)\sin nx \mathrm{d}x $ ，用分部积分法求 $ b_n $ .
然后有 $ \dis f(x) \xlongequal{\textrm{f(x)延拓的部分}} F(x) \xlongequal{\textrm{连续点}} \sum_{n = 1}^\infty b_n\sin nx$
并给出 $ x $ 的区间。 

展成余弦级数同理。



\end{document}
